%% Manuscript for Machine Learning (Springer).
%% Numbers in the text are macros in numbers.tex and related files, generated from the
%% released per-split result files.
\documentclass[pdflatex,sn-basic]{sn-jnl}

\usepackage{graphicx}
\usepackage{multirow}
\usepackage{amsmath,amssymb,amsfonts}
\usepackage{amsthm}
\usepackage[title]{appendix}
\usepackage{xcolor}
\usepackage{textcomp}
\usepackage{booktabs}
\usepackage{tabularx}
\usepackage{algorithm}
\usepackage{algorithmicx}
\usepackage{algpseudocode}
\usepackage{xurl}

\theoremstyle{thmstyleone}
\newtheorem{theorem}{Theorem}
\newtheorem{proposition}{Proposition}
\theoremstyle{thmstylethree}
\newtheorem{definition}{Definition}

\raggedbottom

\newcommand{\hlcl}{\textup{\textsc{HLCL}}}
\newcommand{\fdr}{\mathrm{mFDR}}
\DeclareMathOperator{\E}{\mathbb{E}}
\DeclareMathOperator*{\argmin}{arg\,min}
\DeclareMathOperator{\Bin}{Bin}

\input{numbers}
\input{itemcert}
\input{holm}
\input{fdpratio}
\input{perbaseline}
\input{supp}

\begin{document}

\title[Certified label-level selective prediction]{Certified label-level selective prediction for hierarchical outputs, with an application to LLM annotation}

\author*[1,2]{\fnm{Hongsheng} \sur{Wang}}\email{wanghongsheng@zju.edu.cn}
\author[1]{\fnm{Fei} \sur{Wu}}\email{wufei@zju.edu.cn}

\affil*[1]{\orgdiv{College of Computer Science and Technology}, \orgname{Zhejiang University}, \orgaddress{\street{866 Yuhangtang Road, Xihu District}, \city{Hangzhou}, \postcode{310058}, \state{Zhejiang}, \country{China}}}
\affil[2]{\orgname{Loctek Ergonomic Technology Corp.}, \orgaddress{\street{No. 436 Xueshi Road, Shounan Subdistrict, Yinzhou District}, \city{Ningbo}, \postcode{315100}, \state{Zhejiang}, \country{China}}}

\abstract{%
%ABSTRACT-BEGIN
Selective prediction usually accepts or defers a prediction as a whole. For hierarchical outputs, such as taxonomy paths, the
natural unit of risk is a single label, and the labels of one output are
dependent. We study label-level selective prediction for hierarchical outputs
and propose \hlcl{} (hierarchical label-level certified labeling), which accepts
an ancestor-closed prefix of each predicted path with one threshold per depth,
certifies the thresholds by learn-then-test and deploys the certified threshold
with the lowest modeled review effort. Because the label-level marginal false
discovery rate can be written as one bounded loss per output, \hlcl{}
guarantees, with probability at least {\unboldmath$1-\delta$}, that it does not
exceed a target {\unboldmath$\alpha$} for independent and identically
distributed outputs, without a model of the dependence between labels. As it is also
a weighted average of per-depth rates, a second, exact test exists. We evaluated \hlcl{} on
four public taxonomies of depth two to four annotated by three large language
models (\stat{s@nsettings} settings, 100 random splits each). \hlcl{}
met the target in every setting. Under the same certificate, deciding per label
rather than per item accepted significantly more correct labels in
\stat{s@item@power@better} of \stat{s@nsettings} settings and never fewer, and
in \stat{s@tied@power@better} with a single shared threshold. The modeled review
effort was \stat{s@pb@relmin}\% to \stat{s@pb@relmax}\% lower on average for
\hlcl{} than for each certified baseline that met the target. Against the
strongest such baseline in each setting, it was significantly lower in
\stat{s@wins} settings, higher in \stat{s@losses} and not significantly
different in \stat{s@ties}.
%ABSTRACT-END
}

\keywords{selective prediction, risk control, hierarchical classification, multiple testing, large language models, data annotation}

\maketitle

\section{Introduction}
\label{sec:intro}

Selective prediction lets a model abstain when it is unsure, so that the
predictions it keeps meet a target error rate \citep{chow1970optimum,geifman2017selective}.
Methods with distribution-free guarantees treat a prediction as one unit: it is
kept or deferred as a whole. Many outputs, however, are structured. A path in a
taxonomy is a sequence of labels whose reliability usually drops with depth,
and the labels of one output are dependent because they are predicted from the
same input. For such outputs the natural unit of risk is a single label, and two
questions arise that whole-output methods do not answer: how to bound the error
among the labels that are kept without modeling their dependence, and how to keep
the kept labels consistent with the taxonomy.

Our test case is data annotation. Labels for hierarchical taxonomies are costly
because a person must place each item at several levels, and large language
models (LLMs) have become a common first-pass annotator, since on many
taxonomic tasks they agree with human annotators often enough to be useful
\citep{wang2021gpt3,gilardi2023chatgpt,he2024annollm,tan2024survey}. An LLM is
typically sure about the coarse category of an item and unsure about the fine
one. The selective prediction problem is then: \emph{which of the predicted
labels can be kept without review, with a guarantee on the error of what is
kept, and at what remaining cost?}

Recent work gives this problem statistical guarantees. Conformal labeling
\citep{huang2025selective} applies the Benjamini--Hochberg procedure to conformal
$p$-values and controls the false discovery rate (FDR) of the items whose machine
labels are kept. Probably approximately correct (PAC) labeling
\citep{candes2026pac} chooses an uncertainty threshold such that, with high
probability, the average error over the whole dataset is small, where the labels
supplied by experts count as correct. HyPAC (Zeng et al. 2026, arXiv:2602.02550)
calibrates two thresholds to route between a fast model, a slower model and a
person. Selective classification
\citep{geifman2017selective} bounds the risk of the accepted predictions. All of
these methods decide per \emph{item}. Consider an item whose coarse label is
certain and whose fine label is a coin flip. It is either kept in full, which
admits the error, or reviewed in full, which pays for work that was only needed
one level down. Hierarchical selective classification
\citep{goren2024hierarchical} does back off to a coarser node, but it bounds the
error over all items, including those that back off to the root, rather than the
error among the labels that are kept.

\textbf{Our main claim is that, under one fixed certificate, the unit of
decision changes what can be accepted.} With the same test and the same error
target, deciding per label rather than per item keeps more correct labels, even
with one threshold shared by all depths, and at moderate $\alpha$, when coarse
levels are reliable and fine levels are not, it needs less review. On the Blurb Genre Collection
at $\alpha=0.05$, an item-level rule with exactly our certificate accepts on
average \res{bgc}{0.05}{ItemBest}{power} of the correct labels without review,
even when its statistical test is chosen in its favor. Deciding per label accepts
\res{bgc}{0.05}{HLCL}{power} (Section~\ref{sec:res-unit}).

Moving from items to labels raises two difficulties. The first is statistical.
The labels of one item are predicted from the same text and are therefore
dependent, and a hierarchy adds further dependence between a node and its
descendants. A procedure that treats labels as exchangeable units has no
guarantee for the label-level ratio. Moreover, fitting a high-capacity score on
the data used to set its threshold has been shown empirically to break
selective-risk control in document extraction (Gurram 2026, arXiv:2608.14639). The second difficulty is structural. A
label-level rule that ignores the taxonomy may accept a fine label while rejecting
its parent, and the resulting label set is not a path in the taxonomy.

To address both difficulties, we propose \hlcl{} (hierarchical label-level
certified labeling). \hlcl{} turns repeated LLM samples into a learned score for
every predicted taxonomy node, accepts a prefix of each predicted path, certifies
one threshold per depth on labeled calibration data, and deploys the certified
threshold with the lowest estimated review effort. A reviewer receives the
certified prefix and decides only the remaining levels (Fig.~\ref{fig:overview}).
The method needs only a predicted path and a score per node, so it applies to any
hierarchical classifier; LLM annotation is the application in which we evaluate
it.

\begin{figure}[!tbp]
\centering
\includegraphics[width=\textwidth]{Fig1}
\caption{Overview of \hlcl{}. (a) The rule in an annotation workflow:
machine labels that pass the acceptance rule are written directly, the
remaining levels go to a review queue, and a random audit sample is used for
monitoring and recalibration (Section~\ref{sec:method-system}). (b) Method
steps; steps 1--2 apply to every item, steps 3--4 use the fit split and step 5
the calibration split. (c) An illustrative decision for one item with made-up
scores. Depth 3 clears its own threshold but is not accepted, because its parent
is not (ancestor closure), so the reviewer receives \emph{Medicine} and decides
depths 2 and 3 only}
\label{fig:overview}
\end{figure}

The main contributions of this paper are as follows:
\begin{enumerate}
  \item We formulate \emph{label-level hierarchical selective acceptance}. For
  every predicted node, the task is to decide between accepting and reviewing,
  subject to a high-probability bound on the label-level marginal false
  discovery rate (mFDR) and with an
  explicit model of human effort as the objective
  (Section~\ref{sec:setting}).
  \item We propose \hlcl{}, which adapts learn-then-test calibration to this
  problem. It combines a learned node score built from sample agreement and sibling
  competition, an ancestor-closed acceptance rule with one threshold per depth and
  a calibration in which the candidate thresholds, their order and the test
  statistic are all fixed on data disjoint from the calibration split
  (Section~\ref{sec:method}).
    \item We show that the deployed rule satisfies the label-level bound with
  probability at least $1-\delta$ when items are i.i.d., without any assumption
  on the dependence between labels of the same item, because the label-level
  ratio reduces to one bounded loss per item. Because the same ratio is a
  weighted average of per-depth rates, an exact depth-wise intersection--union
  test of the same null hypothesis is also available. We give the minimum
  evidence that each test needs to certify anything, which explains when each
  one is preferable (Proposition~\ref{prop:evidence}). We also show, at the
  population level and for one threshold, that prefix acceptance dominates item
  acceptance in accepted labels and cost, and we give a sufficient condition
  under which the error target carries over (Proposition~\ref{prop:prefix}).
  Proofs are in Appendix~\ref{app:proofs}.
  \item We evaluate \hlcl{} on four public hierarchical corpora of depth two to
  four, with three LLM annotators from different model families and 100 random
  splits per setting, against six baselines, with ablations under a fixed test
    statistic and a sensitivity analysis. \hlcl{} met the target in every setting,
  and a direct check of the guarantee on a known population confirmed it. In the
  settings in which each certified baseline met the target, the mean modeled
  effort of \hlcl{} was \stat{s@pb@relmin}--\stat{s@pb@relmax}\% lower, with
  significance assessed per setting by paired intervals corrected for
  overlapping splits. Against the strongest certified baseline of each setting,
  most differences are not significant. We report the regimes in which
  label-level acceptance gives little or no gain, and we show that the pooled
  target allows the finest level to run above $\alpha$
  (Sections~\ref{sec:setup}--\ref{sec:discussion}).
\end{enumerate}

The remainder of this paper is organized as follows.
Section~\ref{sec:related} reviews related work. Section~\ref{sec:setting}
formulates the problem, and Section~\ref{sec:method} presents \hlcl{} and its
guarantee. Section~\ref{sec:setup} describes the experimental setup, and
Section~\ref{sec:results} reports the results. Section~\ref{sec:discussion}
discusses theoretical and practical implications and threats to validity, and
Section~\ref{sec:conclusion} concludes.

\section{Related work}
\label{sec:related}

Table~\ref{tab:related} positions \hlcl{} against the closest lines of work
along four dimensions: the unit of decision, whether the taxonomy is respected,
the type of guarantee, and whether human effort is modeled.

\begin{table}[!tbp]
\centering
\footnotesize
\setlength{\tabcolsep}{3pt}
\caption{Representative approaches to deciding which machine labels need human
review. ``Item'' means that an item's labels are accepted or reviewed together.
High-probability guarantees hold with probability at least $1-\delta$ over the
calibration sample (for PAC labeling, over the labeling algorithm on a fixed dataset);
``dataset error'' is the average error over all items (realized for PAC labeling,
population for HyPAC), with reviewed items counted as correct; HyPAC's main
experiments use an asymptotic bound; expected-value guarantees hold for the expected false discovery proportion over calibration and test
draws. ``Budget'' means that a fixed review budget is allocated; ``budget
trade-off'' means that cost and quality are traded off without a fixed budget.
mFDR: marginal false discovery rate (Section~\ref{sec:setting}).}
\label{tab:related}
\begin{tabularx}{\textwidth}{>{\raggedright\arraybackslash}X >{\raggedright\arraybackslash}p{2.1cm} >{\raggedright\arraybackslash}p{1.6cm} >{\raggedright\arraybackslash}p{2.9cm} >{\raggedright\arraybackslash}p{1.5cm}}
\toprule
Approach & Unit of decision & Respects taxonomy & Guarantee & Models human effort \\
\midrule
Uncertainty-guided allocation \citep{li2023coannotating} & item & no & none & budget trade-off \\ \addlinespace[2pt]
Conformal labeling \citep{huang2025selective} & item & no & expected FDR & no \\ \addlinespace[2pt]
PAC labeling \citep{candes2026pac} & item & no & high-prob.\ dataset error & label count or per-source cost \\ \addlinespace[2pt]
HyPAC (Zeng et al. 2026) & item & no & high-prob.\ dataset error & tokens or API price \\ \addlinespace[2pt]
Conformal alignment \citep{gui2024conformalalignment} & item & no & expected FDR & no \\ \addlinespace[2pt]
Trust or escalate \citep{jung2025trust} & item & no & high-prob.\ agreement & no \\ \addlinespace[2pt]
QUORUM routing \citep{quorum2026} & item & no & none on label error & budget \\ \addlinespace[2pt]
Selective classification \citep{geifman2017selective} & item & no & high-prob.\ risk & no \\ \addlinespace[2pt]
Hierarchical selective cls.\ \citep{goren2024hierarchical} & item (back-off) & yes & high-prob.\ accuracy within $\pm\epsilon$ of target & no \\ \addlinespace[2pt]
Per-field risk control (Gurram 2026) & field & no & high-prob.\ per group of fields & no \\ \addlinespace[2pt]
Hierarchical conformal (den Hengst et al. 2025) & prediction set & yes & coverage & no \\ \addlinespace[2pt]
Conformal structured prediction \citep{zhang2025conformal} & prediction set & yes (DAG) & marginal or PAC coverage & no \\ \addlinespace[2pt]
\hlcl{} (this work) & label (node) & yes & high-prob.\ label-level mFDR & yes \\
\bottomrule
\end{tabularx}
\end{table}

\subsection{LLMs as annotators and human--LLM collaboration}
LLMs can replace or assist crowd workers on many text-annotation tasks
\citep{wang2021gpt3,gilardi2023chatgpt}. Prompting strategies improve their
agreement with experts \citep{he2024annollm}, and \citet{tan2024survey} survey
LLM-based annotation and data synthesis.
Because LLM labels are imperfect, several systems keep people in the loop.
CoAnnotating \citep{li2023coannotating} allocates instances between humans and
an LLM by the uncertainty of the LLM output, \citet{wang2024humanllm} train a verifier that scores LLM labels and send
low-scoring labels to human re-annotation, and MEGAnno+
\citep{kim2024meganno} is a collaborative annotation system that lets humans
review LLM labels selectively. More broadly, human-in-the-loop machine learning
\citep{wu2022hitl} combines machine output with human judgment. These systems
decide per item and rely on heuristic uncertainty. \hlcl{} adds a statistical
guarantee on the kept labels and decides per taxonomy level.

The node score of \hlcl{} builds on the observation that agreement among
sampled LLM outputs signals reliability \citep{wang2023selfconsistency,kuhn2023semantic}
and that LLMs can state a confidence, although verbalized confidence tends to be
overconfident \citep{xiong2024llms}. Post hoc calibration such as temperature
scaling \citep{guo2017calibration} improves the reliability of scores. A calibrated score, however, still leaves open where to
cut, and the cut must be justified on labeled data. \hlcl{} does not need a
calibrated score for its guarantee to hold.

\subsection{Selective annotation and routing with guarantees}
\citet{huang2025selective} propose conformal labeling. It computes a conformal
$p$-value per item against the calibration items that the model mislabeled and
then applies the Benjamini--Hochberg procedure \citep{benjamini1995fdr} with a
correction factor $(n+1)/(n_0+1)$, where $n_0$ is the number of mislabeled calibration
items, at target $\alpha$, which controls the expected FDR of the set of items whose
machine labels are used. \citet{gui2024conformalalignment} apply the same
construction to decide when the output of a foundation model can be trusted, and
\citet{jin2023selection} develop the underlying theory of selection with
conformal $p$-values. All three select whole items and assume one label per
item. \hlcl{} differs from them in the unit of decision, in enforcing ancestor
closure and in the form of the guarantee.

\citet{candes2026pac} choose one uncertainty threshold so that, with probability at
least $1-\delta$, the average labeling error over the whole dataset is at most
$\epsilon$, where items sent to experts count as correct. HyPAC (Zeng et al. 2026,
arXiv:2602.02550) calibrates two thresholds to route items between two model
tiers and a person and minimizes annotation cost (tokens or API price) under a
high-probability guarantee on the population error, which is asymptotic with the
bound used in its main experiments. Both are the closest in spirit to our high-probability
bound and cost objective. They differ from \hlcl{} in three ways: they decide per
item, they do not use a taxonomy, and they bound the error over all items rather
than among the kept labels. The two targets are linked but not equivalent. Under the 0-1 loss, and with
reviewed labels counted as correct as in these works, the dataset error equals
the error among kept labels times the share of labels kept, so if half of the
labels are reviewed, the error among the kept labels is twice the dataset error.
Section~\ref{sec:setting-effort} discusses what changes when a reviewer keeps a
wrong accepted prefix.

\citet{jung2025trust} decide when the verdict of an LLM judge can be trusted
and escalate the remaining cases to stronger judges or abstain, with a
high-probability guarantee of agreement with human judgments. \citet{quorum2026} route instances to human
or LLM annotators under a budget using difficulty features, with theory on
posterior concentration and regret but no guarantee on the error of the
resulting labels. \citet{cheng2026multiexpert} use conformal risk control with
abstention for pairwise LLM judging. These methods decide per item.

\subsection{Selective and hierarchical classification}
Selective classification lets a classifier abstain, a trade-off first analyzed
by \citet{chow1970optimum}. \citet{geifman2017selective} give a threshold rule
with a high-probability bound on the risk of the accepted predictions, which is
the direct ancestor of our item-level certified baseline.
\citet{goren2024hierarchical} extend selective classification to hierarchies.
Instead of abstaining, their classifier climbs from its predicted leaf to an
ancestor, and a split-conformal quantile sets one threshold so that, with high
probability, the accuracy of the returned nodes lies within a margin $\epsilon$
of the target. This is the work closest to ours, and our ancestor-closed rule
has the same shape. The differences lie in the target and in the decision.
Hierarchical selective classification bounds the error of the returned node
over all items, including the items that back off to the root, and it chooses
the threshold that returns the most specific nodes. In contrast, we bound the
error among the labels written without review, which is a false discovery rate.
We also certify a separate threshold per depth with two complementary tests and
choose among the certified thresholds by human effort. We include an adaptation
of their method as a baseline (Section~\ref{sec:setup-base}).

Hierarchical text classification learns to predict paths in a taxonomy. Surveys
describe the top-down, per-level decision structure that our rule inherits
\citep{silla2011survey}, and neural models encode the label hierarchy with
attention, graph encoders or external knowledge
\citep{zhou2020hiagm}. These models aim at
accuracy and do not decide which predictions may be stored without review;
\hlcl{} can be placed on top of any of them that outputs a path and a score per
node.

\subsection{Conformal prediction and risk control}
Conformal prediction gives distribution-free guarantees under exchangeability
\citep{vovk2022algorithmic,angelopoulos2023gentle}. \citet{mohri2024language} remove claims from a
generated text until the remainder is certified as factual. This back-off is
close in spirit to our ancestor closure, but its unit is a claim rather than a
taxonomy node, and it has no budget. Gurram (2026, arXiv:2608.14639) decides per
extracted field and controls the selective risk among accepted fields, pooled
or within groups of fields, with guarantees that range from expected risk to
per-group high-probability bounds. Field-level decisions are the granularity we
adopt. That work also documents empirically how score refitting and the
clustering of fields within a document invalidate naive procedures. It treats
fields as flat units and does not accept a coarse unit while routing a fine
one, and its document-level tier bounds a per-document average error, whereas our
tests bound the pooled ratio under i.i.d.\ items. We adapt its field-i.i.d.\ tier
as a baseline. Learn-then-test with per-field thresholds has also been used for
financial documents (Jin et al. 2026, arXiv:2609.20110), and FDR-controlling
calibration for accepting single entities in medical extraction (Shrestha and
Kim 2026, arXiv:2603.00924); both treat fields or entities as flat units.
Hierarchical group-conditional risk control (Salem et al. 2026,
arXiv:2607.24562) also tests every node of a tree with a Bonferroni correction,
but its tree organizes groups of inputs, not the labels of one output. Hierarchical conformal classification
(den Hengst et al. 2025, arXiv:2508.13288; \citealp{mortier2026conformal}),
conformal structured prediction over directed acyclic graphs
\citep{zhang2025conformal} and tree-based conformal prediction for multi-label
data \citep{tyagi2023multilabel} return mixed-level prediction sets with
marginal or PAC coverage guarantees. Coverage requires the truth to lie in
the returned set. We instead require the labels written without review to be
correct, which is a different requirement, and none of these methods models
human effort.

Our calibration is an instance of learn-then-test
\citep{angelopoulos2025learnthen}, which generalizes risk-controlling prediction
sets \citep{bates2021rcps}. Candidate parameters are treated as hypotheses and
tested in a fixed sequence \citep{bauer1991multiple}, which controls the
family-wise error rate without splitting the error budget across candidates. We
use the betting $p$-values of \citet{waudbysmith2024estimating}, which belong to
the game-theoretic approach to anytime-valid inference
\citep{ramdas2023game}, and the intersection--union principle
\citep{berger1982multiparameter}. Conformal risk control
\citep{angelopoulos2024conformal} bounds an expected monotone loss, includes a
hierarchical classification example, and could replace our certificate when an
expectation guarantee suffices. Three elements are new in our setting. The first is a hypothesis family of
per-depth thresholds under ancestor closure. The second is the observation that
the label-level ratio is a weighted average of per-depth rates, which makes an
exact intersection--union test available next to the pooled test, together with
the minimum evidence each test needs and a selection between the two on
held-out data. The third is the effort-based
choice among the certified thresholds.

\subsection{Statistical inference with LLM annotations}
Prediction-powered inference \citep{angelopoulos2023ppi} and its use with LLM
annotations \citep{gligoric2025unconfident} combine many machine labels with a
few human labels to obtain valid confidence intervals for population quantities.
Their goal is an estimate rather than a per-item decision about which labels to
keep. They are therefore complementary to \hlcl{}, which decides which labels
can be kept without review.

\section{Problem formulation}
\label{sec:setting}

\subsection{Hierarchical labels and annotators}

Let $\mathcal{T}$ be a label taxonomy, that is, a rooted tree whose nodes are
labels and whose root-to-leaf paths have length at most $D$. An \emph{item} is a
text $x_i$. Its gold annotation is a path $(y^\star_{i,1},\dots,y^\star_{i,d_i})$
with one node per depth, so every ancestor of a gold node is also gold. This
ancestor closure is the convention of hierarchical corpora such as RCV1
\citep{lewis2004rcv1}, Web of Science \citep{kowsari2017hdltex} and the Blurb
Genre Collection \citep{aly2019hierarchical}.

An annotator maps $x_i$ to a predicted path
$(\hat y_{i,1},\dots,\hat y_{i,\hat d_i})$. In this paper the annotator is an
LLM queried $k$ times with temperature sampling, but the method only needs a
score for every predicted node. A predicted node at depth $d$ is
\emph{correct} if the prediction agrees with the gold path up to depth $d$. A
wrong fine label is therefore an error even when the coarse label is right.

A decision rule with parameters $\lambda$ accepts some of the predicted nodes.
Let $A(\lambda)$ and $W(\lambda)$ denote the number of accepted nodes and the
number of accepted but wrong nodes of a random item. We measure label-level
quality by the marginal false discovery rate (\emph{mFDR}). Unlike the false discovery rate of the multiple-testing
literature, which is an expected proportion, it is a ratio of expectations:
\begin{equation}
  \fdr(\lambda) \;=\; \frac{\E[W(\lambda)]}{\E[A(\lambda)]}.
  \label{eq:fdr}
\end{equation}
This is the long-run share of wrong labels among the labels written without
review. By convention, $\fdr(\lambda)=0$ when $\E[A(\lambda)]=0$, since a rule
that accepts nothing makes no error. The node score is fitted on separate data
(Section~\ref{sec:method-score}) and treated as fixed, so \eqref{eq:fdr} is
defined conditionally on the fitted score. This is the quantity that an operator
monitors, and it does not require the labels of an item to be independent.
\citet{huang2025selective}, in the extended arXiv version of their work, note
that the literature prefers the FDR to the mFDR, because the mFDR cannot be
controlled in expectation under the global null and does not govern the joint
behavior of false and total selections. The high-probability form below
addresses the first point: if every predicted label is wrong, any rule that
accepts something has mFDR one, so it is certified with probability at most
$\delta$. The second point stands, and we state it with the guarantee: our bound
concerns a long-run ratio, not the share of wrong labels in one batch.

\paragraph{Form of the guarantee.}
The rule is chosen from a labeled calibration sample, so $\fdr(\hat\lambda)$ is
a random quantity. We require that
\begin{equation}
  \Pr\bigl(\fdr(\hat\lambda) \le \alpha\bigr) \;\ge\; 1-\delta ,
  \label{eq:guarantee}
\end{equation}
where the probability is taken over the calibration sample. Conformal labeling
\citep{huang2025selective} instead controls the expected false discovery
proportion $\E[\mathrm{FDP}]$ over calibration and test draws. The two
guarantees bound different functionals, and neither implies the other. Ours
fixes the target functional, namely the population ratio \eqref{eq:fdr}, and
guarantees that the \emph{deployed} rule meets it with probability $1-\delta$.
Theirs averages the realized proportion over repetitions. PAC labeling
\citep{candes2026pac} also holds with high probability, but its target is the
realized error on one fixed dataset, with probability over the labeling algorithm and
with expert labels counted as correct. HyPAC (Zeng et al. 2026)
bounds the corresponding population risk, with probability over the calibration
sample. Neither of these targets implies a bound on $\fdr(\lambda)$, and
$\fdr(\lambda)$ does not bound the error of the stored labels in general
(Section~\ref{sec:setting-effort}). Neither ours nor HyPAC bounds the realized
error of one finite batch; PAC labeling does so only for the fixed dataset it
labels. Both
ours and conformal labeling concern long-run rates rather than the share of wrong
labels in one particular batch. The high-probability form suits an
operator who deploys one calibrated rule and must vouch for its error rate. We compare
the two empirically in Section~\ref{sec:results}.

\subsection{Human effort}
\label{sec:setting-effort}


Reviewing an item does not have a fixed cost. A reviewer who is told that the
coarse label is certified and who only chooses the sub-category does less work
than one who annotates the item from scratch. Let $A_{i,d}\in\{0,1\}$ indicate
that the depth-$d$ node of item $i$ was accepted. The cost of item $i$ is
\begin{equation}
  c_i \;=\; \sum_{d=1}^{D} \omega_d\,\mathbf{1}\{A_{i,d}=0\}
  \;+\; c_0\,\mathbf{1}\{\exists d: A_{i,d}=0\},
  \label{eq:effort}
\end{equation}
where $\omega_d\ge0$ with $\sum_d \omega_d=1$ is the share of a full annotation
spent on level $d$, and $c_0$ is a fixed overhead for opening and reading an
item. A fully accepted item costs $0$, and a fully manual one costs $1+c_0$. We
report effort as $\E[c]/(1+c_0)$, the fraction of the cost of annotating the
whole corpus by hand. The overhead matters. Without it, the cheapest rule would
touch nearly every item just to fix its fine level. We use $c_0=0.5$ and equal
$\omega_d$, and we vary $c_0$ in Section~\ref{sec:res-sens}. As in PAC labeling \citep{candes2026pac} and HyPAC, we assume that the
reviewer is correct on the levels left to them; the guarantee
\eqref{eq:guarantee} itself does not need this assumption, because it concerns
accepted labels only. Because correctness is prefix agreement, a reviewer who
keeps the accepted prefix, which is what the cost model charges for, stores a
wrong label at every level below a wrong accepted node. The error rate of all
stored labels can therefore exceed $\fdr(\lambda)$. It is at most
$\Pr\{W(\lambda)\ge1\}\le\E[W(\lambda)]\le\alpha\,\E[A(\lambda)]$ on the
certified event, which is at most $D\alpha$ (Appendix~\ref{app:proofs}). If the
reviewer may overturn accepted labels, the stored error is at most
$\E[W(\lambda)]/D\le\fdr(\lambda)$ for full predicted paths, but the cost model
does not charge for overturning. In our runs the stored error stayed below
$\alpha$ (at most \stat{s@stored@max}$\,\alpha$; Table~\ref{tab:depth}). The parameters $c_0$ and
$\omega_d$ are configurable and should be set from the review interface. The
guarantee does not depend on them, and on the Blurb collection \hlcl{} remains
the cheapest certified method over the range of $c_0$ that we tested
(Section~\ref{sec:res-sens}).

\subsection{Objective}

For a risk level $\alpha$ and a confidence level $1-\delta$, we seek a rule
$\hat\lambda$, chosen from the calibration sample, that satisfies
\eqref{eq:guarantee} and has a low expected cost:
\begin{equation}
  \text{minimize}\;\; \E[c(\hat\lambda)]
  \quad\text{subject to \eqref{eq:guarantee}.}
  \label{eq:objective}
\end{equation}
\hlcl{} selects among certified candidates by estimated cost and carries no
optimality guarantee for \eqref{eq:objective}.
The decision is made per node, so a rule may accept the coarse level of an item
and route only its finer levels to review. The constraint is on the pooled
label-level ratio, which is the error rate of the labels written without review, pooled over levels.
An operator who needs a bound at every level can use the depth-wise test of
Section~\ref{sec:method-cert} alone, because its rejection certifies the
error rate of every depth (Section~\ref{sec:disc-guidelines}).

\begin{definition}[Ancestor closure]
\label{def:closure}
A decision rule is \emph{ancestor-closed} if it accepts a node only when it
accepts all predicted ancestors of that node.
\end{definition}

Closure is required for the output to be usable. A label set that contains a
fine label without its parent is not a path in $\mathcal{T}$. It cannot be
stored in a path-keyed table, and it fails any taxonomy validator. Closure also
matches the task. If the annotator hesitates between \emph{Cardiology} and
\emph{Neurology}, the useful partial output is the certified parent
\emph{Medicine}.

\section{Proposed method: \hlcl{}}
\label{sec:method}

\hlcl{} has four parts: a node score (Section~\ref{sec:method-score}), an
ancestor-closed acceptance rule (Section~\ref{sec:method-rule}), a certificate
(Section~\ref{sec:method-cert}) and the choice of the deployed threshold
(Section~\ref{sec:method-select}). Fig.~\ref{fig:overview} shows how they fit
together, Algorithm~\ref{alg:hlcl} summarizes the procedure,
Section~\ref{sec:method-prefix} relates it to item-level acceptance,
Section~\ref{sec:method-system} covers deployment, and
Section~\ref{sec:method-cost} analyzes its costs.

\subsection{Node score from repeated annotation}
\label{sec:method-score}

The annotator is queried $k$ times with temperature sampling and asked for one
path and a stated confidence per level. At each depth, the $k$ paths that agree
with the prefix chosen so far define a distribution over candidate nodes, and
the \emph{consensus path} follows its mode. A node $v$ at depth $d$ is described
by the following features:
\begin{itemize}
  \item \textbf{vote}: the share of the $k$ samples whose path agrees with the
  consensus up to $v$, which is non-increasing along a path;
  \item \textbf{conditional agreement}: $\mathrm{vote}(v)/\mathrm{vote}(\mathrm{parent}(v))$;
  \item \textbf{sibling margin and entropy}: the gap between the two largest
  conditional shares at depth $d$ and the entropy of the conditional
  distribution;
  \item \textbf{stated confidence}: the mean and standard deviation of the
  confidence that the annotator reports for this level over the supporting
  samples;
  \item the parent's vote, the logarithm of the parent's fan-out in
  $\mathcal{T}$, and a depth indicator.
\end{itemize}
A logistic regression on these features maps each node to a score
$s\in(0,1)$. It is fitted on a \emph{fit split} that is disjoint from the
calibration and test data. Unlike the vote, $s$ is not monotone along a path,
because a child can face weaker competition than its parent. Ancestor closure is
therefore a real constraint on the rule. Keeping the fit split separate prevents
the score-refit leakage described by Gurram (2026).

\subsection{Acceptance rule}
\label{sec:method-rule}

Given thresholds $\lambda=(\lambda_1,\dots,\lambda_D)$, one per depth, the rule
accepts the depth-$d$ node of item $i$ when
\begin{equation}
  A_{i,d} \;=\; \prod_{e\le d}\mathbf{1}\{s_{i,e}\ge\lambda_e\}=1,
  \label{eq:rule}
\end{equation}
where $s_{i,e}$ is the score of the node of item $i$ at depth $e$. A depth with no predicted node has $A_{i,d}=0$. The accepted
set of every item is a prefix of its predicted path and is hence
ancestor-closed (Appendix~\ref{app:proofs}). A reviewer of a partially accepted
item receives the certified prefix and decides only the remaining levels, as
assumed by the cost model \eqref{eq:effort}.

We use one threshold per depth rather than one per node. A per-node rule has as
many parameters as $\mathcal{T}$ has nodes, and it cannot be certified from a
calibration sample of realistic size. A per-depth rule has $D$ parameters and is
unaffected by adding leaves at an existing depth. The candidate set contains the
grid of per-depth threshold vectors built from score quantiles on the fit split.
As a special case, it also contains the vectors with equal thresholds at all
depths.

\subsection{Certificate}
\label{sec:method-cert}

For item $i$, let $A_i(\lambda)=\sum_d A_{i,d}$ and
$W_i(\lambda)=\sum_d A_{i,d}\mathbf{1}\{C_{i,d}=0\}$, where $C_{i,d}$ indicates
that the node is correct. We define the loss
\begin{equation}
  L_i(\lambda) \;=\; W_i(\lambda)-\alpha A_i(\lambda)\;\in\;[-\alpha D,\,(1-\alpha)D].
  \label{eq:loss}
\end{equation}
If the items are i.i.d., so are $L_1(\lambda),\dots,L_n(\lambda)$, and
\begin{equation}
  \E[L(\lambda)]\le 0 \iff \E[W(\lambda)]\le\alpha\,\E[A(\lambda)]
  \iff \fdr(\lambda)\le\alpha
  \label{eq:equiv}
\end{equation}
whenever $\E[A(\lambda)]>0$. A valid $p$-value for the null hypothesis $H_\lambda$ that
$\fdr(\lambda)>\alpha$ is therefore all that is needed. The labels of
one item may be arbitrarily dependent, because each item contributes a single
scalar $L_i$.

\paragraph{Pooled betting test.}
We rescale $L$ to $x=(L+\alpha D)/D\in[0,1]$, so that the rescaled zero is
$t=\alpha$. Under
$H_\lambda$, $\E[x]>t$. The capital
$K_r=\prod_{i\le r}\bigl(1+\ell_i(t-x_i)\bigr)$ with predictable bets
$\ell_i\in[0,\,0.5/(1-t)]$ is then a nonnegative supermartingale, so
$p^{\mathrm{bet}}=\min\{1,1/\max_r K_r\}$ is a valid $p$-value by Ville's inequality
\citep{waudbysmith2024estimating}. We use the predictable plug-in bet
\begin{equation}
  \ell_i=\min\Bigl\{\sqrt{2\log(20)/(n\hat\sigma^2_{i-1})},\;0.5/(1-t)\Bigr\},
  \label{eq:bet}
\end{equation}
where $\hat\sigma^2_{i-1}$ is the running variance of $x_1,\dots,x_{i-1}$,
regularized as $\hat\sigma^2_{j}=\bigl(1/4+\sum_{l\le j}(x_l-\hat\mu_l)^2\bigr)/(j+1)$
with $\hat\mu_l=(1/2+\sum_{u\le l}x_u)/(l+1)$ and $\hat\sigma^2_0=1/4$. The
constant $\log(20)$ is the same for every $\delta$.
The items are visited in a fixed random order. This bet adapts to the small
variance of $L$ at conservative thresholds, but any predictable bet in the
stated range gives a valid $p$-value.

\paragraph{Depth-wise intersection--union test.}
Let $q_d=\E[W_d]/\E[A_d]$ be the error rate among the nodes accepted at depth
$d$. The pooled rate is a weighted average of these rates:
\begin{equation}
  \fdr(\lambda)=\frac{\sum_d\E[W_d]}{\sum_d\E[A_d]}
  =\sum_d\frac{\E[A_d]}{\sum_e\E[A_e]}\,q_d .
  \label{eq:weightedavg}
\end{equation}
Hence $\fdr(\lambda)>\alpha$ implies $q_d>\alpha$ for some $d$, i.e.,
$H_\lambda\subseteq\bigcup_d H_d$ with $H_d:q_d>\alpha$. Because the pairs
$(A_{i,d},W_{i,d})$ are i.i.d.\ across items, the error indicators of the items
accepted at depth $d$ are, conditionally on which items are accepted,
independent with success probability $q_d$. With $n_d$ acceptances and $w_d$
errors,
\begin{equation}
  p_d=\Pr\bigl(\Bin(n_d,\alpha)\le w_d\bigr)
  \label{eq:binomial}
\end{equation}
is an exact $p$-value for $H_d$. If a depth has a finite threshold but no
acceptance in the sample, we set $p_d=1$. Depths with an infinite threshold are
omitted, and the candidate with an infinite threshold at every depth accepts
nothing and is trivially admissible. By the intersection--union principle \citep{berger1982multiparameter},
\begin{equation}
  p^{\mathrm{iut}}=\max_d p_d
  \label{eq:iut}
\end{equation}
is a valid $p$-value for $H_\lambda$ without any multiplicity correction
(Appendix~\ref{app:proofs}). The test depends on the score only through which
nodes are accepted. It is coarser than the pooled test, since it requires every
depth to meet $\alpha$. However, it uses exact tails and remains usable when $\alpha n$ is small. Its rejection also certifies
the stronger statement $q_d\le\alpha$ for every depth that can accept.

\paragraph{Choosing the test.}
Neither test dominates the other. The pooled test exploits the constraint that
is actually the target and is sharper when depths differ in error rate. The
depth-wise test, in contrast, is the only one that certifies anything when
$\alpha$ is very small. \hlcl{} uses the pooled test if the sequence it induces
on the fit split certifies at least one candidate, and the depth-wise test
otherwise. This rule depends only on the fit split, so it does not affect the
guarantee. The following result shows how much evidence each test needs.

\begin{proposition}[Minimum evidence]
\label{prop:evidence}
Let $n_d$ be the number of calibration acceptances at depth $d$ and
$a=\sum_i A_i(\lambda)$ the total number of acceptances. The depth-wise test
can reject $H_\lambda$ at level $\delta$ only if
$n_d\ge\log(1/\delta)/\{-\log(1-\alpha)\}$ at every depth with a finite
threshold. The pooled test with bets in $[0,0.5/(1-\alpha)]$ can reject only if
$a\ge2(1-\alpha)D\log(1/\delta)/\alpha$.
\end{proposition}

The proof is in Appendix~\ref{app:proofs}. Both are necessary, not sufficient,
conditions and assume that no errors are observed; observed errors raise the
requirement. At
$\delta=0.1$ and $D=3$, $\alpha=0.02$ requires $n_d\ge114$ at every active
depth for the depth-wise test or $a\ge677$ for the pooled test, and
$\alpha=0.01$ requires $n_d\ge230$ or $a\ge1{,}368$. The pooled requirement
grows with $D$ and with $1/\alpha$, so at small $\alpha$ the depth-wise test is
the only one that can certify a rule that accepts a few hundred labels. This
matches the transition we observe in Section~\ref{sec:res-abl}.

\begin{theorem}[Validity]
\label{thm:validity}
Assume that the fit, calibration and test items are i.i.d.\ draws, with the
annotator and its sampling fixed, so that the node features of an item depend
only on that item. Let the candidate
set, their order and the choice of test be measurable functions of the fit
split, and let the candidates be tested on the calibration split in that order,
stopping at the first $p$-value above $\delta$. Then, with probability at least
$1-\delta$ over the calibration split and conditionally on the fit split, every
certified candidate $\lambda$ satisfies $\fdr(\lambda)\le\alpha$. In particular,
the deployed $\hat\lambda$ satisfies \eqref{eq:guarantee}.
\end{theorem}

The proof is given in Appendix~\ref{app:proofs}. The guarantee concerns the
deployed rule. Any selection among the certified candidates, including the
effort criterion below, preserves it, because all of them are covered
simultaneously. When the depth-wise test is used, the same argument applies to
the stronger null hypothesis $H'_\lambda:\exists d,\ q_d>\alpha$, so the
certificate then bounds every depth that can accept. Test items are not needed
for the guarantee; in our experiments, the fit, calibration and test parts are
drawn without replacement from a finite pool of 3,000 items.

\subsection{Choosing the deployed threshold}
\label{sec:method-select}

Certification leaves a set of admissible thresholds. Selective annotation
methods usually deploy the most permissive one, which maximizes the number of
accepted labels. \hlcl{} instead estimates the effort \eqref{eq:effort} of each
certified candidate on the calibration split and deploys the cheapest one. The
two criteria differ when a threshold accepts slightly more labels but leaves more
items partially open, since each such item pays the overhead $c_0$.

\begin{algorithm}[t]
\caption{\hlcl{}}
\label{alg:hlcl}
\begin{algorithmic}[1]
\Require items, annotator, taxonomy $\mathcal{T}$, risk $\alpha$, confidence
$\delta$, cost model $(\omega,c_0)$
\State query the annotator $k$ times per item; build consensus paths and node features
\State split the labeled items into fit, calibration and test parts
\State fit the node score $s$ on the fit split
\State build candidate threshold vectors from score quantiles on the fit split
\State choose the test statistic on the fit split (Section~\ref{sec:method-cert})
\State order candidates by their $p$-value on the fit split, conservative first among ties
\For{each candidate $\lambda$ in this order}
  \State compute its $p$-value on the calibration split
  \If{$p>\delta$} \textbf{break}
  \EndIf
  \State record $\lambda$ and its estimated effort $\hat c(\lambda)$
\EndFor
\State \Return the recorded $\lambda$ with the smallest $\hat c$ (accept nothing if none)
\end{algorithmic}
\end{algorithm}

\subsection{Label-level against item-level acceptance}
\label{sec:method-prefix}

The experiments compare the two units of decision under the same certificate. The
following result isolates what can be said without any sampling. For a threshold
$\tau$, let the \emph{item rule} accept all $D$ levels of an item if its predicted path
has $D$ levels and its weakest node has score at least $\tau$, and nothing otherwise.
Let the \emph{prefix rule} be \eqref{eq:rule} with $\lambda_d=\tau$ for all $d$, which
belongs to our candidate set up to the grid. Call an item \emph{partial} if the prefix
rule accepts at least one but fewer than $D$ levels. Let $\Delta A$ and $\Delta W$ be
the accepted and the accepted but wrong nodes of partial items, and let
$\rho=\E[\Delta W]/\E[\Delta A]$ when $\E[\Delta A]>0$.

\begin{proposition}[Prefix against item acceptance]
\label{prop:prefix}
Fix $\tau$ and a score. For every item, the prefix rule accepts every node that the
item rule accepts, and therefore every correct node that it accepts, and
its cost \eqref{eq:effort} is not larger. The expected saving is
\begin{equation}
  \E[c^{\mathrm{item}}]-\E[c^{\mathrm{prefix}}]
  =\E\Bigl[\mathbf{1}\{\text{partial}\}\sum_d\omega_d A^{\mathrm{prefix}}_{d}\Bigr].
  \label{eq:saving}
\end{equation}
If $\E[\Delta A]>0$ and the item rule can accept, then $\fdr(\mathrm{prefix})$ is a
weighted average of $\fdr(\mathrm{item})$ and $\rho$. Hence, if
$\fdr(\mathrm{item})\le\alpha$ and $\rho\le\alpha$, then the prefix rule also
satisfies $\fdr\le\alpha$, accepts at least as many correct nodes and costs no
more. If $\rho>\alpha$, the prefix rule at the same threshold can exceed $\alpha$.
\end{proposition}

The proof is given in Appendix~\ref{app:proofs}. The result is a population
statement about one threshold. It does not say that \hlcl{} beats the item rule on a
given sample, because each rule is certified on its own and neither family
contains the other. The saving \eqref{eq:saving} is the work avoided on partial items and does
not depend on $\rho$. The quantity $\rho$ decides whether the saving is available at the
same threshold: if $\rho\le\alpha$ and the item rule meets the target, it is; if
$\rho>\alpha$, the prefix rule may need a stricter threshold. The saving vanishes when no
item is partial. This is consistent with the regimes of Section~\ref{sec:res-nohelp},
but it does not predict them, because certified rules deploy different thresholds.

\subsection{Deployment}
\label{sec:method-system}

The fitted score, the deployed $\hat\lambda$, the chosen test, $\alpha$, $\delta$ and
the calibration date form a policy that is applied to each new item in one pass over
its levels (Fig.~\ref{fig:overview}(a)). Accepted labels can still be wrong at a
long-run rate of up to $\alpha$, so a reviewer should be able to overturn them.
Reviewer corrections cover only the rejected levels, so they are a selected sample
and cannot be reused as calibration data. Monitoring therefore needs a small random
audit of fully reviewed items, with recalibration when the estimated error drifts or
when the taxonomy, the prompt or the LLM changes.

\subsection{Computational and annotation cost}
\label{sec:method-cost}

Let $n$ be the number of calibration items and $\Lambda$ the candidate set.
Applying rule \eqref{eq:rule} to all items costs $O(nD)$ for one candidate, and
both $p$-values are computed from the per-item losses or per-depth counts in
$O(nD)$ time. Ordering the candidates on the fit split and scanning them on the
calibration split therefore costs $O(|\Lambda|\,nD)$ in the worst case, plus a
sort of $|\Lambda|$ values. With $m$ quantile levels per depth (and an infinite
threshold), the per-depth grid has at most $(m+1)^D$ vectors. We use $m=8$, or
$m=6$ for the four-level Amazon taxonomy, which gives at most 729 vectors for
three levels and 2,401 for four, plus at most 33 vectors with a shared
threshold. Fitting the logistic score is linear in the number of fit-split
nodes. Deployment is a single pass over the new items: computing the scores and
applying the deployed thresholds takes $O(D)$ per item. In practice the cost is
dominated by the $k$ LLM queries per item. In our runs, the 15,000 queries of
one corpus ($3{,}000$ items, $k=5$) took between 15 and 18 minutes of wall-clock
time with \texttt{deepseek-chat}, about 74 minutes with
\texttt{claude-haiku-4-5} and about three and a half hours with
\texttt{gpt-5.6-luna} through a shared, rate-limited gateway, with parallel requests. The calibration itself reads cached responses and needs no API access.
Two further costs matter in deployment. First, the number of samples $k$ is a deployment
parameter that scales the LLM cost linearly; we use $k=5$.
Second, the fit and calibration splits must be labeled by hand. With $n_0$
labeled items and per-item effort $e$ on the remaining items, the effort for a
corpus of $N$ items is $\bigl(n_0+e\,(N-n_0)\bigr)/N$ in units of fully manual
annotation. With $n_0=1{,}800$ and $e=\res{bgc}{0.1}{HLCL}{effort}$ (Blurb
collection, $\alpha=0.1$), this is $0.68$ for $N=3{,}000$, $0.27$ for
$N=20{,}000$ and approaches $e$ for large $N$; the labeled sample is reused for
later batches as long as the item distribution does not change.

\section{Experimental setup}
\label{sec:setup}

The experiments answer five research questions:
\begin{description}
  \item[RQ1] Does the realized label-level error of \hlcl{} respect the target
  at the promised rate (Section~\ref{sec:res-valid})?
  \item[RQ2] How much does deciding per label rather than per item gain under
  the same certificate (Section~\ref{sec:res-unit})?
  \item[RQ3] How much human effort does \hlcl{} save against each baseline,
  and in which regimes does label-level acceptance pay off most
  (Sections~\ref{sec:res-effort} and \ref{sec:res-nohelp})?
  \item[RQ4] Which components drive the gains (Section~\ref{sec:res-abl})?
  \item[RQ5] How sensitive are the results to the review overhead, the
  calibration size, $\delta$ and the annotator
  (Sections~\ref{sec:res-sens} and~\ref{sec:res-annot})?
\end{description}

\subsection{Data and annotators}
\label{sec:setup-data}

We use four public hierarchical corpora whose fine levels differ in difficulty.
\begin{itemize}
    \item \textbf{DBPedia classes} have 9, 70 and 219 labels at three levels.
  They are derived from the DBpedia knowledge base \citep{lehmann2015dbpedia};
  we use the release of \citet{dbpediaclasses2022data} and sample from the first
  11,151 records of its validation split.
  \item \textbf{WOS-46985} \citep{kowsari2017hdltex,kowsari2018wosdata}, a collection of Web of
  Science abstracts, has $7/134$ labels at two levels; we sample from all
  46,985 records.
  \item The \textbf{Blurb Genre Collection} \citep{aly2019hierarchical,bgc2019data} has
  three levels and an official hierarchy with 147 edges. From the English test
  split we keep the blurbs whose genre path is unique and has at least three
  levels.
  \item The \textbf{Home \& Kitchen} product categories of Amazon Reviews 2023
  \citep{hou2026amazon,amazonreviews2023data} have four levels with 8, 31, 112 and 296 nodes. From the
  first 400,000 records of the metadata file, we keep the products whose
  category path has at least four levels below the root, and the paths with at
  least 100 products. This yields a pool of 154,787 products. As text, we use
  the title, the first 500 characters of the features and the first 400 of the
  description, truncated to 1,200 characters in total.
\end{itemize}
From each corpus we draw $3{,}000$ items with a fixed seed. The taxonomy
offered to the annotator is known in advance: the official hierarchy for the
Blurb collection and the label paths of the whole corpus for the other three.
All gold paths and all consensus paths in our samples have full depth. Annotator A
(\texttt{deepseek-chat}) labels all four corpora. Annotators B
(\texttt{claude-haiku-4-5}) and C (\texttt{gpt-5.6-luna}), which come from two
further model families, label the Blurb collection again. Each item is queried $k=5$ times at temperature $1.0$ with the prompt given in
Appendix~\ref{app:prompt}. All 90,000 responses are cached and released, so every number below can
be reproduced without API access. Four responses could not be parsed: three
blank ones on WOS-46985 and one truncated one from annotator B. They count as
missing samples. On Amazon, three responses left every level empty, and they
count as votes for no label.

Table~\ref{tab:annot} gives the per-depth accuracy of the consensus path.
DBPedia is easy at every level. The Blurb collection degrades at the third
level for all three annotators: annotator B is less accurate than annotator A at
every depth, and annotator C is slightly more accurate. Amazon degrades steadily with depth, from
$0.97$ at the first level to $0.73$ at the fourth. WOS-46985 is hard at the
second level, where the annotator is right less than half the time. In that
regime, no certified method can accept much at small $\alpha$.

\begin{table}[!tbp]
\centering
\caption{Accuracy of the consensus path by depth. Annotator A is
\texttt{deepseek-chat}, annotator B is \texttt{claude-haiku-4-5} and annotator C
is \texttt{gpt-5.6-luna}.}
\label{tab:annot}
\input{annot}
\end{table}

\subsection{Baselines}
\label{sec:setup-base}

All methods use the same learned node score $s$ and the same splits.
\begin{itemize}
  \item \textbf{Item-level certified selection} (Item-cert.). Each item is
  scored by its weakest node, and a single threshold is certified, so the item
  is accepted or reviewed as a whole. This is the high-probability selective rule of \citet{geifman2017selective}
  applied to items, with our certificate and 64 candidate thresholds (score
  quantiles of the fit split).
  We run it with the automatic test choice of \hlcl{}, with the pooled test only
  and with the depth-wise test only. In every table we report the variant with
  the lowest mean effort in that setting, denoted Item-cert.$^\ast$. Choosing
  this variant on the test results is an oracle choice that can only favor the
  baseline.
  \item \textbf{Hierarchical selective classification} (HSC)
  \citep{goren2024hierarchical}. Each item returns the deepest prefix of its
  predicted path whose nodes all clear one shared threshold. The original method
    scores a node by the aggregated probability of its descendants under a flat
  classifier. We use our node score instead, so that all methods share one
  score. Following that work, the risk is the item-level error over all items,
  where an item backed off to the root counts as correct. The original sets the
  threshold to an order statistic of per-item minimal thresholds, with a
  two-sided guarantee around the target. To give a one-sided certificate
  comparable with the other baselines, we instead certify with an exact binomial
  tail and fixed-sequence testing over 64 candidate thresholds and deploy the
  most permissive certified one, which keeps the original preference for the
  most specific nodes. Its guarantee concerns a different quantity from ours. We
  report it because it is a recent published hierarchical back-off method.
  \item \textbf{Per-field selective risk control} (Per-field)
    (Gurram 2026, arXiv:2608.14639), adapted. Each depth is one group of fields
  with one threshold, certified with an exact binomial tail at level $\delta/D$,
  which is the simultaneous form of that work's field-i.i.d.\ tier. The original
  tests 15 quantile candidates with a mix of Holm and fixed-sequence testing; we
  use 200 fit-split score quantiles in a fixed sequence ordered on the fit split,
  so that the order is independent of the calibration data. The most permissive
  certified threshold is deployed, with ties broken towards the most permissive
  candidate, and closure is not enforced.
  \item \textbf{Item-level conformal selection} (Item-conf.)
  \citep{huang2025selective,jin2023selection}. It computes a conformal $p$-value
    per item from the item score and applies Benjamini--Hochberg at target
  $\alpha$ with the $(n+1)/(n_0+1)$ correction. It controls the
  \emph{expected} item-level FDR. It selects only items with a full predicted
  path, so the item-level FDP bounds the label-level FDP pointwise, but it does
  not give \eqref{eq:guarantee}. It uses the same fitted score but does not use the fit split for ordering or
  test selection.
  \item \textbf{Depth-wise Benjamini--Hochberg} (Depth-BH). It computes
  conformal $p$-values per node and applies Benjamini--Hochberg separately at
  each depth. The nodes of one depth are i.i.d.\ across items, so this method
  controls the expected FDP at each depth. It gives no high-probability
  statement and no direct statement about the pooled ratio.
  \item \textbf{Fixed threshold} ($s\ge0.9$). This is the common practitioner
  rule, and it carries no guarantee.
\end{itemize}
Item-cert., HSC and Per-field carry high-probability certificates, and we call
them the \emph{certified baselines}. Item-conf.\ and Depth-BH carry
expected-value guarantees. HSC and Per-field deploy the most permissive
certified threshold, as in their original form, whereas \hlcl{} deploys the
cheapest one; with the most permissive criterion, the effort of \hlcl{} changes
by at most a few points (Section~\ref{sec:res-abl}). A method counts as
\emph{on target} in a setting if its realized label-level error exceeded
$\alpha$ in at most $\delta=10\%$ of the splits. This is an operational
criterion, not a test of \eqref{eq:guarantee}, which concerns the population
ratio; Section~\ref{sec:res-valid} adds a direct check. A method that is not on
target in a setting accepts too many wrong labels for the label-level
requirement, which also lowers its effort, so its effort is not comparable. We
therefore compare \hlcl{} with each baseline over the settings in which that
baseline is on target. In all \stat{s@inval@HSC} settings excluded for HSC, it was cheaper than
\hlcl{}, so Section~\ref{sec:res-effort} also reports the comparison over all
settings. The classification changes little with the cut-off, except for
item-level conformal selection: with 20\% instead of $\delta$, HSC, item-level
conformal selection, depth-wise Benjamini--Hochberg and the fixed threshold are
not on target in \stat{s@inval20@HSC},
\stat{s@inval20@ConfLabel}, \stat{s@inval20@DepthBH} and
\stat{s@inval20@Conf>=0.9} settings, against \stat{s@inval@HSC},
\stat{s@inval@ConfLabel}, \stat{s@inval@DepthBH} and \stat{s@inval@Conf>=0.9}
with $\delta$. As a stricter test, we also compare
\hlcl{} in each setting with the \emph{strongest certified baseline}, the
certified baseline on target with the lowest mean effort in that setting. This baseline is
selected with knowledge of the test results, which favors the baselines.

\subsection{Protocol and metrics}
\label{sec:setup-proto}

For each corpus, annotator and $\alpha$, we draw 100 random splits into fit,
calibration and test parts of 600, 1200 and 1200 items, respectively. The score,
the candidates, their order and the test statistic are fixed on the fit split.
Certification uses the calibration split only. Unless stated otherwise, all
reported numbers are computed on the test split with $\delta=0.1$. We report
three metrics: (i) \emph{power}, the share of correct label nodes accepted
without review; (ii) \emph{effort} \eqref{eq:effort}, as a fraction of fully
manual annotation; and (iii) \emph{viol.}, the share of splits whose realized
label-level error on the test split exceeds $\alpha$. The guarantee concerns the
population rate, so viol.\ is a proxy: a finite test split can exceed $\alpha$
when the population rate does not, and the reverse, and the 100 splits are not
independent replications. We therefore also check the guarantee directly
(Section~\ref{sec:res-valid}). Effort is the \emph{modeled} review effort
\eqref{eq:effort}; we did not time reviewers. The values of $\alpha$ differ by
corpus, because a useful range depends on the annotator's accuracy.

To compare two methods in a setting, we use the paired difference of their
effort over the 100 splits. The splits are drawn from the same 3,000 items, so
the differences are positively correlated and the naive standard error is too
small. We therefore use the corrected resampled-$t$ variance of
\citet{nadeau2003inference}. It multiplies the sample variance by
$1/J+n_{\mathrm{test}}/n_{\mathrm{train}}=\stat{s@nbfactor}$, with $J=100$
splits, $n_{\mathrm{test}}=1200$ and $n_{\mathrm{train}}=1800$. \hlcl{}
\emph{wins} a setting if the corrected 95\% interval for \hlcl{} $-$ baseline
lies below zero. It \emph{loses} if the interval lies above zero, and it \emph{ties} otherwise.
The interval uses the normal quantile 1.96; with the $t_{99}$ quantile (1.984),
two wins, both on WOS-46985 (against Item-cert.$^\ast$ at $\alpha=0.2$ and
Depth-BH at $\alpha=0.3$), would become ties, and no other classification changes. We write paired differences as the mean $\pm$ the
half-width of this interval. Because \stat{s@nsettings} settings are compared, we also report
how many wins and losses survive a Holm adjustment over the settings. Across
settings, we summarize the certified methods by their average ranks, the
Friedman test and the Nemenyi critical difference \citep{demsar2006statistical}.
Because the values of $\alpha$ of one corpus--annotator pair share items and
annotations, we repeat the Friedman test with these pairs as the units; the
three annotators of the Blurb collection still share items, so even these units
are not fully independent.

\subsection{Implementation details}
\label{sec:setup-impl}

The node score is an $\ell_2$-regularized logistic regression (inverse
regularization strength 1.0) on features standardized on the fit split. The
values of $\alpha$ were set per corpus from the consensus accuracy
(Table~\ref{tab:annot}), which was computed on all 3,000 items, so that each
range spans small to moderate risk. The method and the grids of $\alpha$ were
developed iteratively on the same corpora with annotators A and B, and the grid
of the Blurb collection was changed once, from
$\{0.05,0.1,0.15,0.2\}$ to $\{0.02,0.05,0.1,0.15\}$, after early runs. The
candidate order and the test choice use the in-sample scores of the fit split,
which we preferred to a cross-fitted variant in a pilot; both are functions of
the fit split only and do not affect validity. Annotator C was added after the code
was frozen, so its four settings are a check on data not used in development.
Candidate thresholds use $m=8$ score quantiles per depth on the fit split ($m=6$ on Amazon), plus an infinite threshold that
switches a depth off, and $4m$ quantiles for the shared-threshold candidates.
The pooled betting test uses the plug-in bet of Section~\ref{sec:method-cert};
the depth-wise test uses exact binomial tails. All methods are implemented in
Python with NumPy and scikit-learn. The code, the cached LLM responses, the
per-split result files and the scripts that regenerate every table and figure
are publicly available (see the Declarations). Rerunning the experiments on a
different machine can change single splits through floating-point ties among
candidate $p$-values; in our reruns of all 24 settings, 55 of 2,400 splits of
\hlcl{} changed, mean effort changed by at most $0.008$ (DBPedia at
$\alpha=0.02$) and by at most $0.002$ elsewhere, and the significance counts of
Section~\ref{sec:res-unit} were unchanged.

Apart from the annotators under study, we used Claude (Anthropic) as a coding and
writing assistant for the analysis and plotting scripts and for drafting and
editing parts of the text. The authors reviewed the code and the text, and all
results were produced by running the released code.

\section{Results}
\label{sec:results}

\begin{table}[!tbp]
\centering
\footnotesize
\setlength{\tabcolsep}{2.5pt}
\caption{Effort/power over 100 splits; lower effort is better. The first four
method columns carry high-probability certificates. Item-conf.\ and Depth-BH
control an expectation, and $s\ge0.9$ has no guarantee. Item-cert.$^\ast$ is
the best of three test choices in each row. Bold marks the lowest effort among
the certified methods that met the label-level target in the row. $^\dagger$
marks a realized error above $\alpha$ in more than $10\%$ of splits.}
\label{tab:main}
\input{main_results}
\end{table}

\begin{figure}[!tbp]
\centering
\includegraphics[width=\textwidth]{Fig2}
\caption{Power (top row; share of correct labels accepted without review) and
human effort (bottom row; fraction of fully manual annotation) against the risk
level $\alpha$ for \hlcl{} and the three certified baselines. Each point is the mean over 100 splits; A, B and C denote the three annotators.
Hollow markers show settings in which the method's realized label-level error
exceeded $\alpha$ in more than 10\% of splits, so its effort there is not
comparable}
\label{fig:power}
\end{figure}

\begin{figure}[!tbp]
\centering
\includegraphics[width=\textwidth]{Fig3}
\caption{Share of the 100 splits whose realized label-level error on the test
split exceeded $\alpha$, for all methods. The dashed line marks $\delta=0.1$; a
method with a high-probability guarantee should stay near or below it}
\label{fig:validity}
\end{figure}

Table~\ref{tab:main} and Figs.~\ref{fig:power}--\ref{fig:validity} summarize
the results. The paired and rank-based comparisons follow in
Sections~\ref{sec:res-unit} and~\ref{sec:res-effort}.

\subsection{Validity (RQ1)}
\label{sec:res-valid}
In every setting, the realized error of \hlcl{} exceeded $\alpha$ in at most
\stat{s@maxviol@HLCL} of splits, which is consistent with $\delta=0.1$
(Fig.~\ref{fig:validity}). \hlcl{} was thus on target in all \stat{s@nsettings}
settings. Because viol.\ is only a proxy, we also checked
\eqref{eq:guarantee} directly. We treated the 3,000 items of each setting as the
population, drew fit and calibration samples from them with replacement 100
times, and computed the exact mFDR of each deployed rule on all 3,000 items. The
share of draws in which it exceeded $\alpha$ was at most
\stat{s@pop@HLCL@max} for \hlcl{}, \stat{s@pop@ItemLTT@max} for Item-cert.\ and
\stat{s@pop@PerField-LTT@max} for Per-field, all within $\delta$
(Table~\ref{tab:popcheck} in Appendix~\ref{app:supp}). For HSC, whose
certificate concerns the item-level error over all items, the label-level mFDR
exceeded $\alpha$ in up to \stat{s@pop@HSC@max} of draws, on WOS-46985, in line
with the test-split results below. On average, the realized error of \hlcl{} was between
\stat{s@fdpratio@HLCL@min}$\,\alpha$ and \stat{s@fdpratio@HLCL@max}$\,\alpha$,
so the certificate keeps a safety margin below the target. The methods
reach different realized errors at the same nominal $\alpha$ (median
\stat{s@ratio@HLCL@med}$\,\alpha$ for \hlcl{} against
\stat{s@ratio@ItemBest@med}$\,\alpha$, \stat{s@ratio@HSC@med}$\,\alpha$ and
\stat{s@ratio@PerField-LTT@med}$\,\alpha$ for the three certified baselines;
Table~\ref{tab:realized}), so part of the lower effort of \hlcl{} reflects a
less conservative use of the same error budget. All effort comparisons below
are at equal nominal $\alpha$, not at equal realized error. Item-cert.\ and Per-field exceeded $\alpha$ in at most
\stat{s@maxviol@ItemBest} and \stat{s@maxviol@PerField-LTT} of splits,
respectively.
HSC certifies the item-level error over all items, which is a different
quantity. When the annotator is weak, its label-level error among the accepted
labels can be far above $\alpha$. On WOS-46985 at $\alpha=0.1$, it exceeded
$\alpha$ in \res{wos}{0.1}{HSC}{viol} of splits. The reason is that over a third
of the items back off to the root and count as correct, which leaves room for
errors among the accepted labels. We therefore treat HSC as a competitor
only in the rows where it met the label-level target. Depth-wise
Benjamini--Hochberg exceeded $\alpha$ in up to \stat{s@maxviol@DepthBH} of
splits, although its mean realized error stays close to $\alpha$, at most
\stat{s@maxfdpratio@DepthBH}$\,\alpha$. Depth-BH controls the expected error at each depth, which does not bound the
pooled ratio, so these exceedances are not a failure of that method. The fixed threshold
exceeded $\alpha$ in up to \stat{s@maxviol@Conf>=0.9} of splits.
Item-level conformal selection exceeded $\alpha$ in up to
\stat{s@maxviol@ConfLabel} of splits.

In total, HSC, Item-conf., Depth-BH and the fixed threshold, which certify or
control other quantities or nothing, were not on target in
\stat{s@inval@HSC}, \stat{s@inval@ConfLabel}, \stat{s@inval@DepthBH} and
\stat{s@inval@Conf>=0.9} of the \stat{s@nsettings} settings, respectively.

\paragraph{Error by depth.}
The guarantee bounds the error pooled over levels, and coarse labels are easier,
so the pooled ratio can hide a higher error at the finest level.
Table~\ref{tab:depth} reports the error among the accepted labels of each depth.
At the finest level it reached up to \stat{s@depth@finalmax}$\,\alpha$, for
example about $0.17$ on the Blurb collection at $\alpha=0.1$, while the first
level stayed below $\alpha$ in every setting (at most $0.85\,\alpha$). An operator who needs the target at every
level must use the depth-wise test, whose rejection bounds every depth
(Section~\ref{sec:method-cert}).

\subsection{Unit of decision (RQ2)}
\label{sec:res-unit}
Item-cert.\ carries exactly the certificate of \hlcl{}. Since we report its best
test choice per setting, the comparison isolates the unit of decision, together with its per-depth
thresholds, and, if anything, favors the baseline. Table~\ref{tab:itemcert} gives the corrected paired differences in power and
effort for all \stat{s@nsettings} settings. Deciding per label accepted significantly more
correct labels in \stat{s@item@power@better} settings (\stat{s@holm@item@power@better}
after the Holm adjustment) and significantly fewer in none. It needed
significantly less effort in \stat{s@item@effort@better} settings
(\stat{s@holm@item@effort@better} after adjustment) and significantly more in
none. The gains are largest at moderate risk. On the Blurb collection at $\alpha=0.05$, \hlcl{} accepts
\res{bgc}{0.05}{HLCL}{power} against \res{bgc}{0.05}{ItemBest}{power} for
annotator A, \res{bgc_haiku}{0.05}{HLCL}{power} against
\res{bgc_haiku}{0.05}{ItemBest}{power} for annotator B and
\res{bgc_gpt}{0.05}{HLCL}{power} against \res{bgc_gpt}{0.05}{ItemBest}{power} for
annotator C. On WOS-46985 at
$\alpha=0.2$, it accepts \res{wos}{0.2}{HLCL}{power} against
\res{wos}{0.2}{ItemBest}{power}. The gap shrinks as $\alpha$ approaches the annotator's own error rate, where
both rules accept almost every label, which is consistent with \eqref{eq:saving}, since few items are then partial.

To separate the unit of decision from the per-depth thresholds, we compare the
tied variant of Section~\ref{sec:res-abl}, which decides per label with one
threshold shared by all depths and the automatic test, with Item-cert.$^\ast$
(Table~\ref{tab:tied} in Appendix~\ref{app:supp}). It accepted significantly more
correct labels in \stat{s@tied@power@better} settings
(\stat{s@tied@power@holm} after the Holm adjustment) and significantly fewer in
\stat{s@tied@power@worse}, and it needed significantly less effort in
\stat{s@tied@effort@better} settings (\stat{s@tied@effort@holm} after
adjustment) and more in \stat{s@tied@effort@worse}. The unit of decision alone
thus accounts for most of the gain in accepted labels. At the smallest $\alpha$
of DBPedia, the Blurb collection and Amazon, the tied variant and Item-cert.$^\ast$
both accept almost nothing, and the gains of \hlcl{} there come from per-depth
thresholds. Accepted labels are not equally valuable. On the Blurb collection at
$\alpha\ge0.1$, all of the additional correct labels lie at the coarser levels,
and the item-level rule accepts up to $0.06$ more of the finest-level labels,
whereas at $\alpha=0.05$ \hlcl{} also accepts more finest-level labels, for
example $0.47$ against $0.25$ of them with annotator A (\hlcl{} values in
Table~\ref{tab:depth}).

\begin{table}[!tbp]
\centering
\footnotesize
\caption{Paired differences \hlcl{} $-$ Item-cert.$^\ast$ (same certificate,
decision per item) in power and effort over 100 splits: mean $\pm$ half-width
of the Nadeau--Bengio-corrected 95\% interval. Positive power and negative
effort differences favor \hlcl{}; bold marks significant differences in favor
of \hlcl{}. No difference is significant in favor of Item-cert.$^\ast$.}
\label{tab:itemcert}
\input{itemcert_table}
\end{table}

\subsection{Review effort against the baselines (RQ3)}
\label{sec:res-effort}

\begin{table}[!tbp]
\centering
\footnotesize
\setlength{\tabcolsep}{2.5pt}
\caption{Effort of \hlcl{} against each baseline, over the settings in which the
baseline is on target (realized error above $\alpha$ in at most 10\% of splits).
Wins and losses are significant paired differences (Nadeau--Bengio-corrected
95\% intervals over 100 splits); Holm w/l counts those that remain after a Holm
adjustment over the settings. Reduction is the relative difference of the mean
efforts.}
\label{tab:perbaseline}
\input{perbaseline_table}
\end{table}

Table~\ref{tab:perbaseline} compares \hlcl{} with each baseline over the
settings in which that baseline is on target. Against the three certified
baselines, \hlcl{} needed on average \stat{s@pb@item@rel}\%,
\stat{s@pb@hsc@rel}\% and \stat{s@pb@pf@rel}\% less review effort than
Item-cert.$^\ast$, HSC and Per-field, respectively. It was significantly cheaper
in \stat{s@pb@item@wins}, \stat{s@pb@hsc@wins} and \stat{s@pb@pf@wins} settings
and significantly more expensive in \stat{s@pb@item@losses},
\stat{s@pb@hsc@losses} and \stat{s@pb@pf@losses}; after the Holm adjustment,
\stat{s@pb@item@holmwins}, \stat{s@pb@hsc@holmwins} and
\stat{s@pb@pf@holmwins} of these wins remain. These are ratios of mean
efforts over the selected settings; the median per-setting reduction is
\stat{s@medred@ItemBest}\%, \stat{s@medred@HSC}\% and
\stat{s@medred@PerField-LTT}\%, respectively. On the Blurb collection at $\alpha=0.05$, for example, the
effort falls from \res{bgc}{0.05}{PerField-LTT}{effort} (Per-field) and
\res{bgc}{0.05}{HSC}{effort} (HSC) to \res{bgc}{0.05}{HLCL}{effort}. Against the
item-level rule with the same certificate, the gap narrows as $\alpha$ grows; on
Amazon at $\alpha=0.1$, \hlcl{} needs \res{amazon}{0.1}{HLCL}{effort} and
Item-cert.$^\ast$ \res{amazon}{0.1}{ItemBest}{effort}. Over all settings,
including those in which a baseline missed the target, the reductions are
\stat{s@pb@item@allrel}\%, \stat{s@pb@hsc@allrel}\% and
\stat{s@pb@pf@allrel}\%, and HSC is significantly cheaper in
\stat{s@pb@hsc@alllosses} settings, \stat{s@pb@hsc@losses} of which is a setting
in which it met the target. Against the expected-value baselines and the fixed
threshold, the reductions are larger in the settings in which these methods met
the target.

\begin{figure}[!tbp]
\centering
\includegraphics[width=\textwidth]{Fig4}
\caption{Average ranks of the certified methods by effort over the
\stat{s@fr@N} settings (lower is better; a method that is not on target in a
setting is ranked last there). Methods joined by a bar do not differ by more than the
Nemenyi critical difference (CD) at the 5\% level. The settings of one
corpus--annotator pair share items, so the diagram is descriptive
(Section~\ref{sec:res-effort})}
\label{fig:cd}
\end{figure}

Across settings and ranked by effort alone, \hlcl{} has the best average rank
among the four certified methods, \stat{s@fr@erank@HLCL}, followed by
\stat{s@fr@enext@name} with \stat{s@fr@enext}; this gap \stat{s@fr@enext@sig}
the critical difference. When a method that misses the target in a setting is
ranked last there (Fig.~\ref{fig:cd}), the average ranks are
\stat{s@fr@rank@HLCL} for \hlcl{}, \stat{s@fr@rank@ItemBest} for
Item-cert.$^\ast$, \stat{s@fr@rank@HSC} for HSC and
\stat{s@fr@rank@PerField-LTT} for Per-field (Friedman test, $p\stat{s@fr@p}$).
With the \stat{s@fr@bN} corpus--annotator pairs as units and their mean effort,
the average ranks are \stat{s@fr@brank@HLCL}, \stat{s@fr@brank@ItemBest},
\stat{s@fr@brank@HSC} and \stat{s@fr@brank@PerField-LTT} (Friedman test,
$p\stat{s@fr@bp}$); with six units the critical difference is 1.91, and only the
gap to Per-field exceeds it. We read the rank analysis as a summary that
supports the per-setting intervals, because the settings of one pair are not
independent.

\begin{figure}[!tbp]
\centering
\includegraphics{Fig5}
\caption{Paired effort difference between \hlcl{} and the on-target certified
baseline with the lowest mean effort in each setting (marker shape: which
baseline is strongest in that setting), with
Nadeau--Bengio-corrected 95\% intervals over 100 splits. Negative values favor
\hlcl{}; the line color marks significant wins, significant losses and
non-significant differences}
\label{fig:forest}
\end{figure}

The strongest certified baseline of each setting is a stricter reference,
because it is chosen with knowledge of the test results (Fig.~\ref{fig:forest}).
\hlcl{} has the lower mean effort in \stat{s@lowermean} of the
\stat{s@nsettings} settings. It is significantly cheaper in \stat{s@wins}
settings, significantly more expensive in \stat{s@losses} and not significantly
different in the remaining \stat{s@ties}; after the Holm adjustment,
\stat{s@holm@best@wins} wins and \stat{s@holm@best@losses} loss remain. Against
a baseline selected per setting in this way, most differences are therefore not
significant. Where this baseline is HSC, \hlcl{} additionally provides the
label-level guarantee that HSC lacks; where it is Item-cert.$^\ast$ or
Per-field, it provides the same or a comparable certificate together with an
ancestor-closed output.

\subsection{Operating regimes (RQ3)}
\label{sec:res-nohelp}
The results show where the advantage of label-level acceptance is largest and
where it narrows. It is largest at moderate $\alpha$, when the coarse levels
are reliable and the fine levels are not: on the Blurb collection at
$\alpha=0.05$ and $0.1$ with annotators A and C and on Amazon at $\alpha=0.05$,
\hlcl{} is significantly cheaper even than the strongest certified baseline of
the setting. It narrows in two regimes. First, at the
smallest $\alpha$ of DBPedia, the Blurb collection and Amazon, the calibration
sample contains few expected errors ($\alpha\,n_{\mathrm{cal}}\le24$), so any
certificate on the label-level ratio needs more data. \hlcl{} then uses the
depth-wise test, which requires every depth to meet $\alpha$. Where HSC is on target, it is cheaper
there, significantly so only on Amazon at $\alpha=0.02$, because it certifies a
different target, the item-level error over all items, which is easier to meet
when almost every item has a correct coarse label. Second, when $\alpha$
approaches the error rate of the finest level, almost every label can be
accepted, and deciding per item is as cheap as deciding per label, as on DBPedia
at $\alpha\ge0.05$ and at $\alpha=0.15$ on Amazon and with annotator C, where the item-level rule is
slightly, though not significantly, cheaper. In total, a certified baseline on target
had a lower mean effort than \hlcl{} in \stat{s@cheaper@cert} of the
\stat{s@nsettings} settings, and the difference was significant in
\stat{s@losses}. When the annotator is weak at the fine level, as on WOS-46985,
every certified method routes most items to review at small $\alpha$.
Section~\ref{sec:disc-guidelines} turns these regimes into a configuration
guide.

\subsection{Ablations (RQ4)}
\label{sec:res-abl}

\begin{table}[!tbp]
\centering
\footnotesize
\setlength{\tabcolsep}{2pt}
\caption{Effort of \hlcl{} variants over 100 splits; lower is better. Bold marks
values within $0.005$ of the best in the row and test group. Each variant is run
with the automatic choice of test (left), with the pooled betting test only
(middle) and with the depth-wise test only (right). This design prevents the
effect of a component from being confounded with the test that happens to be
selected. \emph{vote}: sample agreement (with the stated confidence as a tie-breaker)
instead of the learned score;
\emph{no cl.}: no ancestor closure, with nodes accepted independently;
\emph{tied}: one threshold for all depths. With Nadeau--Bengio-corrected
95\% intervals, the full method and the \emph{vote} or \emph{no cl.} variant
differ significantly in no setting; the \emph{tied} variant differs
significantly in \stat{s@abl@tied@auto@sig} settings with the automatic test,
\stat{s@abl@tied@wsr@sig} with the pooled test and \stat{s@abl@tied@iut@sig}
with the depth-wise test.}
\label{tab:abl}
\input{ablations}
\end{table}

Table~\ref{tab:abl} varies one component at a time under a fixed test.

\paragraph{Test statistic.}
The choice of test has by far the largest effect. With the pooled test only,
\hlcl{} certifies almost nothing at the smallest $\alpha$ of DBPedia and of the
Blurb collection with annotators A and B. Its effort is
\res{dbpedia}{0.01}{HLCL-wsr}{effort} on DBPedia at $\alpha=0.01$ and
\res{bgc}{0.02}{HLCL-wsr}{effort} on the Blurb collection at $\alpha=0.02$. With
the depth-wise test only, it is much more expensive at larger $\alpha$; on the
Blurb collection at $\alpha=0.1$, its effort is \res{bgc}{0.1}{HLCL-iut}{effort}
against \res{bgc}{0.1}{HLCL-wsr}{effort}. The selection rule picked the
depth-wise test in all or nearly all splits at the smallest $\alpha$ of every
corpus, and in none of the splits at the two largest values. Near the transition the choice is less clear-cut, and guideline 2 in
Section~\ref{sec:disc-guidelines} gives a rule based on $\alpha\,n_{\mathrm{cal}}$.

\paragraph{Score, closure and per-depth thresholds.}
Per-depth thresholds matter most with the depth-wise test, where a shared
threshold must be strict enough for the hardest depth. For example, the tied variant needs
\res{bgc_haiku}{0.05}{HLCL-tied-iut}{effort} against
\res{bgc_haiku}{0.05}{HLCL-iut}{effort} with annotator B at $\alpha=0.05$, and
\res{amazon}{0.05}{HLCL-tied-iut}{effort} against
\res{amazon}{0.05}{HLCL-iut}{effort} on Amazon at $\alpha=0.05$. With the pooled
test, per-depth thresholds mostly matter at small $\alpha$; on DBPedia at
$\alpha=0.02$, the efforts are \res{dbpedia}{0.02}{HLCL-tied-wsr}{effort} and
\res{dbpedia}{0.02}{HLCL-wsr}{effort}, a difference that is not significant.
The learned score and the raw vote give similar effort. The score is
numerically cheaper on DBPedia at $\alpha=0.02$
(\res{dbpedia}{0.02}{HLCL-wsr}{effort} against
\res{dbpedia}{0.02}{HLCL-vote-wsr}{effort} with the pooled test) and at the
larger $\alpha$ of WOS-46985 and annotator B, the vote on Amazon and with
annotator C, but none of
these paired differences is significant. The score is therefore best treated as
a plug-in component selected on held-out data
(Section~\ref{sec:disc-guidelines}). Closure is a requirement on the output
rather than an efficiency device: without it, the accepted labels need not form
a path in the taxonomy (Definition~\ref{def:closure}). Under the pooled test,
closure never increased effort by more than $0.01$; on DBPedia at
$\alpha=0.02$ the closed rule was cheaper, \res{dbpedia}{0.02}{HLCL-wsr}{effort}
against \res{dbpedia}{0.02}{HLCL-flat-wsr}{effort}, though not significantly.
Under the depth-wise test, the non-closed variant is up to six points cheaper,
again not significantly, and this saving is overstated: its output is not a
path, and the cost model does not charge for the accepted children that a reviewer must re-check after correcting a rejected parent. The same applies to
Per-field, which does not enforce closure. The gains of \hlcl{} thus come mainly
from the label-level unit of decision at moderate $\alpha$, where one shared
threshold is not significantly different from the full grid under the pooled
test, and from per-depth thresholds with the depth-wise test at the smallest
$\alpha$. The learned score gave no significant gain over the raw vote. Closure
is a requirement on the output, not a source of savings.

\paragraph{Objective.}
Deploying the certified candidate with the most accepted labels, instead of the
cheapest one, changes effort by at most a few points at $c_0=0.5$. This variant
is not shown in Table~\ref{tab:abl}. The largest change occurs on Amazon at
$\alpha=0.05$, with \res{amazon}{0.05}{HLCL-maxlab}{effort} against
\res{amazon}{0.05}{HLCL}{effort}.

\subsection{Sensitivity (RQ5)}
\label{sec:res-sens}

\begin{table}[!tbp]
\centering
\footnotesize
\setlength{\tabcolsep}{3pt}
\caption{Sensitivity of \hlcl{} to the review overhead $c_0$, the calibration
size $n_{\mathrm{cal}}$ and the confidence level $\delta$ on the Blurb
collection with annotator A over 100 splits. ``Fully auto'' is the share of
items accepted at every level. The last three columns give the effort of the
certified baselines in the same setting. Here Item-cert.\ uses the automatic
test selection, not the per-setting best variant of Table~\ref{tab:main}.}
\label{tab:sens}
\input{sens}
\end{table}

Table~\ref{tab:sens} varies one design parameter at a time on the Blurb
collection.

\paragraph{Review overhead.}
The certificate does not depend on $c_0$; only the choice among the certified
thresholds does. A larger overhead pushes the deployed rule towards accepting
whole items. At $\alpha=0.05$, the share of fully automatic items rises from
\stat{s@sens@bgc_c0_0@0.05@fullauto} at $c_0=0$ to
\stat{s@sens@bgc_c0_1@0.05@fullauto} at $c_0=1$, while power changes from
\stat{s@sens@bgc_c0_0@0.05@power} to \stat{s@sens@bgc_c0_1@0.05@power}. The
violation share is nearly unchanged, and \hlcl{} has the lowest mean effort among the
certified methods for every value of $c_0$.

\paragraph{Cost weights.}
The equal weights $\omega_d$ assume that every level costs the same to decide,
whereas a fine level with hundreds of options plausibly costs more. We evaluated
the deployed rules of \hlcl{} and Item-cert.$^\ast$ under fine-heavy weights
$\omega_d\propto d$ and under $c_0=2$ on the \stat{s@extra@nsettings} settings of
Table~\ref{tab:costsens} in Appendix~\ref{app:supp}. The number of settings in
which \hlcl{} is significantly cheaper falls from \stat{s@cost@effort@wins} to
\stat{s@cost@eff_lin@wins} and \stat{s@cost@eff_c2@wins}, respectively, and it
is significantly more expensive in none. In particular, the significant wins on
the Blurb collection at $\alpha=0.1$ do not survive either change.

\paragraph{Calibration size.}
For these runs we used the first $n_{\mathrm{cal}}$ items, in corpus order, of
each randomly drawn calibration split. As expected, the
certificate becomes more conservative with fewer calibration items. At $\alpha=0.05$, effort rises from \stat{s@sens@bgc@0.05@effort} with
1,200 items to \stat{s@sens@bgc_ncal600@0.05@effort} with 600 and
\stat{s@sens@bgc_ncal300@0.05@effort} with 300. At $\alpha=0.1$, the item-level
rule in its automatic variant loses more. Its effort rises from $0.28$ to
$0.49$, whereas that of \hlcl{} rises from $0.20$ to $0.28$. With only 300
calibration items and $\alpha=0.05$, HSC, which certifies the item-level error
over all items, is cheaper than \hlcl{}. The violation share stays small
throughout.

\paragraph{Confidence level.}
Changing $\delta$ from $0.05$ to $0.2$ moves effort by about three points at both values of $\alpha$. The violation share tracks $\delta$ as it should. At
$\alpha=0.1$, it is \stat{s@sens@bgc_delta0.05@0.1@viol},
\stat{s@sens@bgc@0.1@viol} and \stat{s@sens@bgc_delta0.2@0.1@viol} of splits,
respectively.

\subsection{Annotators of different accuracy (RQ5)}
\label{sec:res-annot}

On the Blurb collection, annotator B is less accurate than annotator A at every
depth and annotator C is slightly more accurate (Table~\ref{tab:annot}). The
score and the procedure were applied to all three annotators without any
change, and \hlcl{} was on target at every $\alpha$ for each of them: the realized
error exceeded $\alpha$ in at most \res{bgc_gpt}{0.02}{HLCL}{viol} of splits.
With annotators A and C, whose first two levels are accurate, \hlcl{} was
significantly cheaper than the strongest certified baseline at $\alpha=0.05$
and $0.1$; with annotator C, its effort was \res{bgc_gpt}{0.05}{HLCL}{effort}
against \res{bgc_gpt}{0.05}{HSC}{effort} for HSC at $\alpha=0.05$ and
\res{bgc_gpt}{0.1}{HLCL}{effort} against \res{bgc_gpt}{0.1}{ItemBest}{effort}
for Item-cert.$^\ast$ at $\alpha=0.1$. With the weaker annotator B, \hlcl{} has
the lowest mean effort among the certified methods on target at all four values of
$\alpha$, by one to three points, but none of these differences is significant
after the Nadeau--Bengio correction (Fig.~\ref{fig:forest}). Annotator C was
added after the code was frozen, and its results agree with those of annotator
A. The advantage of
label-level acceptance thus depends on the annotator through the accuracy of
its coarse levels, which guideline 1 in Section~\ref{sec:disc-guidelines} asks
to measure first.

\section{Discussion}
\label{sec:discussion}

\subsection{Theoretical implications}
\label{sec:disc-theory}

Four findings bear on how selective acceptance of machine labels should be modeled.
First, under the same certificate, the unit of decision changes what can be
accepted, as Proposition~\ref{prop:prefix} shows at the population level for one
threshold. Deciding per label rather than per item accepted significantly more
correct labels in \stat{s@item@power@better} of \stat{s@nsettings} settings and
never fewer, and in \stat{s@tied@power@better} settings with a single shared
threshold (Section~\ref{sec:res-unit}). Second, because the label-level
ratio is a weighted average of per-depth rates, a pooled test and an exact
depth-wise test are both valid tests of the same null hypothesis. The
depth-wise test certifies a stronger statement, and neither dominates in effort;
which one is cheaper depends mainly on the expected number of errors in the
calibration sample, $\alpha\,n_{\mathrm{cal}}$, and on how much the depths
differ in error rate (Section~\ref{sec:res-abl}), as the minimum evidence of
Proposition~\ref{prop:evidence} predicts. Third, a pooled label-level target is weaker than it looks: it allowed the
finest level to run up to \stat{s@depth@finalmax}$\,\alpha$
(Section~\ref{sec:res-valid}), and it bounds the error of the stored labels only by $D\alpha$ when a reviewer
keeps a wrong prefix (Section~\ref{sec:setting-effort}). Fourth, the two forms of
guarantee answer different needs. The high-probability certificates rarely
exceeded $\alpha$ in more than $\delta$ of the splits, whereas item-level
conformal selection and depth-wise Benjamini--Hochberg did so in
\stat{s@inval@ConfLabel} and \stat{s@inval@DepthBH} settings. This is expected by
construction: viol.\ is a quantile criterion, and the latter two methods control
expectations. In the settings in which they were on target, \hlcl{} needed on average \stat{s@pb@conf@rel}\% and
\stat{s@pb@dbh@rel}\% less effort, partly because item-level conformal
selection controls the stricter item-level error. An expected-value guarantee
suits an operator who recalibrates often and cares about average quality over
many deployments. An operator who deploys one calibrated rule and must vouch for
its error rate needs the high-probability form. Both forms can share the node
score and the ancestor-closed rule.

\subsection{Guidance for practitioners}
\label{sec:disc-guidelines}


Table~\ref{tab:guide} summarizes which rule to use in which regime, based on
Tables~\ref{tab:main}, \ref{tab:abl} and~\ref{tab:sens}. The following
guidelines add the operating details.

\begin{table}[!tbp]
\centering
\footnotesize
\setlength{\tabcolsep}{3pt}
\caption{Which acceptance rule to use, based on our experiments and the certificate of Section~\ref{sec:method-cert}.
$\alpha\,n_{\mathrm{cal}}$ is the target error rate times the number of
calibration items.}
\label{tab:guide}
\begin{tabularx}{\textwidth}{>{\raggedright\arraybackslash}p{4.2cm} >{\raggedright\arraybackslash}X >{\raggedright\arraybackslash}p{3.3cm}}
\toprule
Regime & Recommended rule & Evidence \\
\midrule
Coarse levels reliable, fine levels not; moderate $\alpha$ & \hlcl{} with the pooled test & Blurb (A, C) at $\alpha=0.05$--$0.1$, Amazon at $\alpha=0.05$; Tables~\ref{tab:main} and~\ref{tab:abl}, Fig.~\ref{fig:forest} \\ \addlinespace[2pt]
$\alpha$ close to the error rate of the finest level & \hlcl{} or the item-level rule with the same certificate; both accept almost every label & DBPedia at $\alpha\ge0.05$, Table~\ref{tab:main} \\ \addlinespace[2pt]
Very small $\alpha$ ($\alpha\,n_{\mathrm{cal}}\le24$) & \hlcl{} with the depth-wise test, or more calibration data; HSC if an item-level accuracy target over all items suffices & smallest $\alpha$ on DBPedia, Blurb and Amazon, Table~\ref{tab:abl} \\ \addlinespace[2pt]
A bound is needed at every level & \hlcl{} with the depth-wise test only & Section~\ref{sec:method-cert} \\ \addlinespace[2pt]
Few labeled items ($n_{\mathrm{cal}}\approx300$) & \hlcl{}; expect a more conservative certificate; at small $\alpha$, HSC can be cheaper if an item-level accuracy target suffices & Table~\ref{tab:sens} \\ \addlinespace[2pt]
No labeled data & label a calibration sample first; a fixed threshold gives no guarantee & Fig.~\ref{fig:validity} \\
\bottomrule
\end{tabularx}
\end{table}

\begin{enumerate}
  \item \emph{Measure the annotator first.} Estimate the per-depth accuracy of
  the consensus path on a small labeled sample (Table~\ref{tab:annot}).
  Label-level acceptance pays off most when the coarse levels are reliable and
  the fine levels are not. When $\alpha$ approaches the annotator's error rate,
  an item-level rule is as cheap; when $\alpha$ is far below the accuracy of
  even the coarsest level, little can be certified by any method.
  \item \emph{Choose the test from $\alpha\,n_{\mathrm{cal}}$.} With 1,200
    calibration items, the depth-wise test was cheaper in five of the six
  settings with $\alpha\,n_{\mathrm{cal}}\le24$, in line with
  Proposition~\ref{prop:evidence}, and the pooled test was cheaper
  or at most 0.012 more expensive in every setting with $\alpha\,n_{\mathrm{cal}}\ge60$
  (Table~\ref{tab:abl}); with 300 and 600 calibration items, the pooled test
  still certified rules at $\alpha\,n_{\mathrm{cal}}=30$ (Table~\ref{tab:sens}).
  The automatic rule of Section~\ref{sec:method-cert} follows this pattern
  except near the transition and on WOS-46985 at $\alpha=0.1$, where both tests
  certify almost nothing, and the chosen test should be logged with the
  policy.
    \item \emph{Decide whether a pooled bound suffices.} The certificate bounds the
  error pooled over levels, and in our runs the finest level ran up to
  \stat{s@depth@finalmax}$\,\alpha$. If a downstream use relies on the finest
  level, certify with the depth-wise test, whose rejection bounds the error at
  every level.
  \item \emph{Budget the labeled data.} Our runs used 600 fit and 1,200
  calibration items. Section~\ref{sec:method-cost} shows how this one-off cost
  is amortized: for the Blurb collection at $\alpha=0.1$, total effort falls
  from 0.68 of fully manual annotation for 3,000 items to 0.27 for 20,000 items.
    \item \emph{Treat the score as a replaceable component.} The learned score and the raw
  vote did not differ significantly in our runs (Section~\ref{sec:res-abl}).
  Both are valid inputs to the certificate. Choosing between them on held-out
  data is a reasonable rule, but we did not test it.
  \item \emph{Calibrate the cost model.} Set $c_0$ and $\omega_d$ from the review
    process. The certificate does not depend on them, so a wrong guess costs
  efficiency, not validity; the size of the saving, however, does depend on them
  (Section~\ref{sec:res-sens}).
\end{enumerate}

\subsection{Threats to validity}
\label{sec:disc-limits}
\emph{Distribution shift.} As for any calibration-based method, the guarantee
assumes that calibration and deployment items come from the same distribution.
Items that share a source should be treated at the level of the independent
unit (Gurram 2026). The random audit of Section~\ref{sec:method-system}
can detect drift and trigger recalibration; we did not evaluate this loop.

\emph{Scope of the guarantee.} The guarantee bounds the long-run error rate of
the labels accepted without review, pooled over levels. It does not bound the
finest level, which ran up to \stat{s@depth@finalmax}$\,\alpha$, nor the error
of one batch, nor that of the stored labels below a wrong accepted node. When
every level must meet the target, the depth-wise test provides a bound per
level.

\emph{Statistical conclusion validity.} The 100 splits reuse the same 3,000
items per corpus, and the corrected resampled-$t$ interval is an approximation
for this design. The values of $\alpha$ of one corpus--annotator pair share
items and annotations, so the settings are not independent, and we treat the
rank tests across settings as a summary.

\emph{Construct validity.} Correctness is defined by the gold paths of the
corpora, some of which (Amazon) are assigned by sellers and noisy. The corpora
are public, so the annotators may have seen them in training; this can affect the
measured savings but not the validity of the certificate, which is computed on
labeled data. The values of $\alpha$ were chosen per corpus from the consensus
accuracy of all 3,000 items, and the method was developed on the same data,
except for annotator C (Section~\ref{sec:setup-impl}). Effort is modeled, not
measured. The cost model assumes a correct reviewer and linear savings per
accepted level; significant wins over the item-level rule at $\alpha=0.1$ on the
Blurb collection do not survive fine-heavy weights or $c_0=2$
(Section~\ref{sec:res-sens}).

\emph{Baseline configuration.} HSC and Per-field are adaptations that share our
node score and deploy the most permissive certified threshold, as in their
original form, whereas \hlcl{} deploys the cheapest one. Deploying the most
permissive threshold changes the effort of \hlcl{} by at most a few points
(Section~\ref{sec:res-abl}), so the baselines keep their original
deployment rule.

\emph{External validity.} We evaluated four public corpora of depth two to four,
three annotators from different model families, one prompt and $k=5$ samples.
All predicted and gold paths in our samples have full depth, so the case of
shorter paths, which the method allows, is not exercised. Taxonomies with
several gold paths per item or a directed acyclic structure require a different
notion of correctness.

\section{Conclusions and future work}
\label{sec:conclusion}

We studied selective prediction at the level of single labels of a
hierarchical output, under a certified error bound and an explicit reviewing
cost, and evaluated it on LLM annotation. \hlcl{}
keeps the accepted labels ancestor-closed, certifies per-depth thresholds with
a learn-then-test procedure whose design choices are all fixed on held-out
data, and deploys the cheapest certified threshold. On four public corpora of
depth two to four with three LLM annotators, \hlcl{} met the target in all
\stat{s@nsettings} settings, whereas \stat{s@inval@nbase} of the 6 baselines
missed it in up to \stat{s@inval@max} settings. In the settings in which each
certified baseline met the target, the mean modeled effort of \hlcl{} was
\stat{s@pb@relmin}--\stat{s@pb@relmax}\% lower, and it had the best average rank
among the certified methods. Against the strongest baseline chosen per setting
with knowledge of the test results, it was significantly cheaper in
\stat{s@wins} settings and costlier in \stat{s@losses}, so its advantage is
concentrated where coarse levels are reliable and fine levels are not, at
moderate $\alpha$. Under the same certificate, deciding
per label rather than per item accepted significantly more correct labels in
\stat{s@item@power@better} of \stat{s@nsettings} settings and never fewer, for
example \res{bgc}{0.05}{HLCL}{power} instead of \res{bgc}{0.05}{ItemBest}{power}
on the Blurb Genre Collection at $\alpha=0.05$. We recommend measuring the
per-depth accuracy of the annotator first, choosing the test from the
calibration budget, and using the depth-wise test when the finest level must meet
the target on its own.

Future work offers several opportunities: adaptive test selection that uses the
calibration size, sharper per-level and batch-level guarantees, guarantees under
distribution shift between calibration and deployment, label-level procedures
with expected-value guarantees, multi-path and DAG taxonomies, and a timed user
study of review effort in a real annotation interface.

\backmatter

\section*{Statements and Declarations}

\bmhead{Funding}
No external funding was received for conducting this study. The employer of the
first author provided access to an LLM gateway (see Competing interests).

\bmhead{Competing interests}
Hongsheng Wang is employed by Loctek Ergonomic Technology Corp., whose e-mail
annotation workload motivated this study. The company provided access to the LLM
gateway through which annotator C was queried. It had no role in the design of the
method, the choice of datasets, the analysis or the decision to publish, and no
company data are used. Fei Wu declares no competing interests.

\bmhead{Data and code availability}
The four corpora are public: DBPedia classes \citep{dbpediaclasses2022data},
released under CC0; WOS-46985 \citep{kowsari2018wosdata}; the Blurb Genre
Collection \citep{bgc2019data}, released under CC BY-NC; and the Home \&
Kitchen metadata of Amazon Reviews 2023 \citep{amazonreviews2023data}. We redistribute no text
from any corpus; users obtain the corpora from their sources under the
providers' terms. The code, the cached LLM responses (labels and stated
confidences only, no source texts), the per-split result files and the scripts
that regenerate every table and figure are available at
\url{https://github.com/wanghongsheng01/hlcl} and archived at Zenodo
(\url{https://doi.org/10.5281/zenodo.23762584}).

\bmhead{Ethics approval and consent}
Not applicable.

\bmhead{Authors' contributions}
Hongsheng Wang: conceptualization, methodology, software, validation, formal
analysis, investigation, data curation, writing (original draft, review and
editing), visualization. Fei Wu: supervision, writing (review and editing).

\begin{appendices}

\section{Proofs}
\label{app:proofs}

\paragraph{Ancestor closure of the acceptance rule.}
If $A_{i,d}=1$, then $\mathbf{1}\{s_{i,e}\ge\lambda_e\}=1$ for every $e\le d$,
and hence $A_{i,e}=1$ for every $e\le d$.

\paragraph{Validity of the depth-wise $p$-value.}
Fix $\lambda$ and $d$. The pairs $(A_{i,d},W_{i,d})$ are i.i.d.\ across items,
with $W_{i,d}\le A_{i,d}$. Conditionally on $(A_{1,d},\dots,A_{n,d})$, the
variables $W_{i,d}$ of the items with $A_{i,d}=1$ are independent Bernoulli
variables with parameter $\Pr(W_{i,d}=1\mid A_{i,d}=1)=q_d$. Hence
$W_d=\sum_i W_{i,d}\mid n_d\sim\Bin(n_d,q_d)$. Under $H_d:q_d>\alpha$, the
binomial tail \eqref{eq:binomial} satisfies $\Pr(p_d\le t\mid n_d)\le t$ for all
$t$, because the $\Bin(n_d,q)$ distribution is stochastically increasing in $q$.
Integrating over $n_d$ gives $\Pr(p_d\le t)\le t$. If $n_d=0$ and
$\lambda_d<\infty$, we set $p_d=1$, which is trivially valid. If
$\lambda_d=\infty$, the rule cannot accept at depth $d$, so $\E[A_d]=0$. Depth
$d$ then does not enter \eqref{eq:weightedavg} and is omitted from the maximum.
By \eqref{eq:weightedavg}, $\fdr(\lambda)>\alpha$ implies that $H_d$ holds for
some $d$, say $d^\ast$. Then $\Pr(\max_d p_d\le t)\le\Pr(p_{d^\ast}\le t)\le t$,
so $p^{\mathrm{iut}}$ is valid for $H_\lambda$. If the test rejects, every
$H_d$ with $\E[A_d]>0$ is rejected at level $\delta$, so the same certificate
bounds $q_d$ at every depth that can accept.

\paragraph{Proof of Proposition~\ref{prop:prefix}.}
If the item rule accepts an item, every node has score at least $\tau$, so the prefix
rule accepts all of them too. This gives the pointwise inclusion of accepted nodes and
of correct accepted nodes. Items fall into three classes. If the prefix rule accepts
all $D$ levels, both costs are $0$. If it accepts none, both costs are $1+c_0$. If the item is
partial, the item rule accepts nothing and costs $\sum_d\omega_d+c_0$, whereas the prefix rule costs
$\sum_d\omega_d\mathbf{1}\{A_d=0\}+c_0$, so the difference is
$\sum_d\omega_d A_d$. Taking expectations gives \eqref{eq:saving}. The accepted nodes of
the prefix rule are those of the item rule plus those of partial items, so
$\E[W^{\mathrm{prefix}}]=\E[W^{\mathrm{item}}]+\E[\Delta W]$ and likewise for $A$.
Writing $a=\E[A^{\mathrm{item}}]$ and $b=\E[\Delta A]$, we obtain
$\fdr(\mathrm{prefix})=\bigl(a\,\fdr(\mathrm{item})+b\,\rho\bigr)/(a+b)$, a weighted
average. The claim for $\rho\le\alpha$ follows. The last claim needs an example. Take $D=2$ and
$\alpha=0.1$. With probability $0.1$ both levels are accepted and correct. With probability
$0.9$ only level 1 is accepted, and it is wrong with probability $0.2$. Then
$\fdr(\mathrm{item})=0$, $\rho=0.2$ and $\fdr(\mathrm{prefix})=0.18/1.1>\alpha$. In general,
the prefix rule exceeds $\alpha$ if and only if $b(\rho-\alpha)>a(\alpha-\fdr(\mathrm{item}))$.

\paragraph{Proof of Proposition~\ref{prop:evidence}.}
Depth-wise test: $p_d=\Pr\{\Bin(n_d,\alpha)\le w_d\}\ge\Pr\{\Bin(n_d,\alpha)=0\}
=(1-\alpha)^{n_d}$, and $p^{\mathrm{iut}}\le\delta$ requires $p_d\le\delta$ at
every depth with a finite threshold, hence $(1-\alpha)^{n_d}\le\delta$. Pooled
test: with $t=\alpha$ and $x_i=(W_i-\alpha A_i+\alpha D)/D$, we have
$t-x_i=(\alpha A_i-W_i)/D\le\alpha A_i/D$. Since $\log(1+u)\le u$ and
$0\le\ell_i\le0.5/(1-\alpha)$,
$\log K_r\le\sum_{i\le r}\ell_i\,\alpha A_i/D\le
0.5\,\alpha\,a/\{(1-\alpha)D\}$ for every $r$. Rejection requires
$\max_r K_r\ge1/\delta$, hence $a\ge2(1-\alpha)D\log(1/\delta)/\alpha$.

\paragraph{Error of the stored labels.}
Let a reviewer keep the accepted prefix of an item and complete its path. If no
accepted node is wrong, the stored path is correct by assumption. Otherwise at
most $D$ stored labels are wrong. For full-depth paths, the expected share of
wrong stored labels is therefore at most $\Pr\{W(\lambda)\ge1\}\le\E[W(\lambda)]$, which is at most
$\alpha\,\E[A(\lambda)]\le D\alpha$ when $\fdr(\lambda)\le\alpha$. The stored error can
exceed $\fdr(\lambda)$: take $D=2$, let half of the items have both levels
accepted and correct and half have only level 1 accepted, wrong with
probability $0.2$. Then $\fdr=0.1/1.5\approx0.067$, but $10\%$ of the stored
labels are wrong. If the reviewer may also overturn accepted labels, only
accepted wrong labels can remain, so the share is at most $\E[W(\lambda)]/D$.

\paragraph{Proof of Theorem~\ref{thm:validity}.}
Condition on the fit split. The candidates $\lambda_{(1)},\lambda_{(2)},\dots$,
their order and the test statistic are then fixed. The calibration items are
i.i.d.\ and independent of them, so each $p_{(j)}$ is a valid $p$-value for
$H_{\lambda_{(j)}}$. Fixed-sequence testing rejects $H_{\lambda_{(1)}},\dots,
H_{\lambda_{(J)}}$, where $J$ is the index before the first $p$-value above
$\delta$. Let $j^\ast$ be the smallest index whose null hypothesis is true. A
true null is rejected only if $J\ge j^\ast$, which requires
$p_{(j^\ast)}\le\delta$. Hence the probability of rejecting any true null is at
most $\Pr(p_{(j^\ast)}\le\delta)\le\delta$ \citep{bauer1991multiple}. On the
complementary event, which has probability at least $1-\delta$, every certified
$\lambda$ has $\fdr(\lambda)\le\alpha$ by \eqref{eq:equiv}. A candidate with
$\E[A(\lambda)]=0$ accepts nothing and is harmless. The deployed $\hat\lambda$
is one of the certified candidates or accepts nothing, which proves
\eqref{eq:guarantee}. The argument holds conditionally on the fit split and
therefore also unconditionally.

\section{Annotation prompt}
\label{app:prompt}

Every item was annotated with the following prompt, sent as a single user
message with temperature 1.0, at most 160 output tokens and JSON output mode
where the API supports it. \texttt{\{taxonomy\}} lists the taxonomy with one line per deepest parent, in
the form ``\texttt{a > b > leaf1 | leaf2 | ...}''; \texttt{\{text\}} is the item
text truncated to 1,500 characters. The prompt is shown for a three-level
taxonomy; for two and four levels, the header and the JSON keys have two or four
levels. The full taxonomy is included in every prompt.

\begin{footnotesize}
\begin{verbatim}
Classify the text into exactly one path of the taxonomy below
(coarse to fine). Use label names exactly as written.

TAXONOMY (each line: level-1 > level-2 > level-3 options):
{taxonomy}

TEXT:
{text}

Answer with JSON only: {"l1", "l2", "l3", "c1", "c2", "c3"} where c* is
your probability (0-1) that the label at that level is correct.
\end{verbatim}
\end{footnotesize}

\section{Supplementary results}
\label{app:supp}

All tables in this appendix use the splits of Table~\ref{tab:main}. Tables
\ref{tab:depth}, \ref{tab:popcheck} and~\ref{tab:costsens} come from a rerun of the deployed rules
from the cached responses on the machine used for this revision, which can differ
from the released runs in single splits (Section~\ref{sec:setup-impl}).

\begin{table}[!tbp]
\centering
\footnotesize
\caption{Error and power of \hlcl{} by depth, mean over 100 splits. $q_d/\alpha$
is the error among the labels accepted at depth $d$ divided by $\alpha$ (--: no
acceptance at that depth); power is the share of correct labels of each depth
that are accepted; stored is the error of all stored labels when a reviewer keeps
the accepted prefix and completes the path, divided by $\alpha$. $q_d$ is
averaged over the splits with at least one acceptance at depth $d$ and is noisy
where power is near zero.}
\label{tab:depth}
\input{depth_table}
\end{table}

\begin{table}[!tbp]
\centering
\footnotesize
\caption{Direct check of the guarantee. The 3,000 items of each setting are
treated as the population; in each of 100 draws, fit and calibration samples are
drawn from them with replacement, and the exact mFDR of the deployed rule is
computed on all 3,000 items. The table gives the share of draws in which it
exceeds $\alpha$; for \hlcl{}, Item-cert.\ (automatic test choice) and
Per-field, the guarantee requires at most $\delta=10\%$.}
\label{tab:popcheck}
\input{popcheck_table}
\end{table}

\begin{table}[!tbp]
\centering
\footnotesize
\setlength{\tabcolsep}{3pt}
\caption{Mean realized label-level error divided by $\alpha$ over 100 splits.
All methods are compared at the same nominal $\alpha$, but they use the error
budget to different degrees.}
\label{tab:realized}
\input{realized_table}
\end{table}

\begin{table}[!tbp]
\centering
\footnotesize
\caption{Paired differences between the tied variant, which decides per label
with one threshold shared by all depths, and Item-cert.$^\ast$, with
Nadeau--Bengio-corrected 95\% intervals over 100 splits. Bold marks significant
differences in favor of the tied variant.}
\label{tab:tied}
\input{tied_table}
\end{table}

\begin{table}[!tbp]
\centering
\footnotesize
\caption{Effort difference \hlcl{} $-$ Item-cert.$^\ast$ under the default cost
model, under fine-heavy weights $\omega_d\propto d$ and under $c_0=2$, for the
same deployed rules, with Nadeau--Bengio-corrected 95\% intervals. Bold marks
significant differences in favor of \hlcl{}.}
\label{tab:costsens}
\input{costsens_table}
\end{table}

\end{appendices}

\bibliography{references}

\end{document}
